\chapter{初识 SLAM}
\label{ch:intro}

% ════════════════════════════════════════════════════════════════
%  全书开篇章：吸收 视觉SLAM十四讲 ch2「初识SLAM」(主骨架) + Cadena2016「Past,
%  Present and Future of SLAM」 + Durrant-Whyte & Bailey 2006 Part I + SLAM
%  Handbook(Visual SLAM / Factor Graphs)。自包含、全吸收，是后续每章的地图。
%  教学层遵循 v5.0 理论教学规范（动机→反面→历史→理论→陷阱→练习 +
%  符号表/速查/故障排查），与黄金范本 lie_theory.tex 同结构。
% ════════════════════════════════════════════════════════════════

\begin{practice}[前置自测]
开始之前，先试答下列问题。它们勾勒了本章——也是全书——要回答的核心疑问。若已能流畅作答，可只速读本章的框架图、路线图与速查表；若卡住，正是本章要为你打通的。
\begin{enumerate}
  \item 一台机器人要在陌生房间里自由活动，至少必须回答哪两个问题？这两个问题之间为什么会构成“鸡生蛋”的循环？
  \item 为什么不能像装 GPS、铺导轨那样“一劳永逸”地解决机器人定位？“未知环境”这四个字为何是 SLAM 的命门？
  \item 一个经典视觉 SLAM 系统由哪几个模块组成？把相邻两帧图像的运动一段段“串”起来，能不能直接得到一条可信的长轨迹？为什么不能？
  \item 把“机器人带着传感器在未知环境里运动”翻译成数学，会得到哪两个方程？方程里的噪声项意味着 SLAM 本质上是什么类型的问题？
  \item 求解这个问题，历史上有滤波与优化两条路线。今天主流为何倒向了优化（图优化）？滤波又被边缘化到了哪些场景？
\end{enumerate}
\end{practice}

\section*{本章目标}
学完本章，你应当能够：
\begin{enumerate}
  \item \textbf{说清} SLAM 是什么、为什么需要它，以及定位与建图为何必须\emph{联合}估计（“鸡生蛋”问题的由来与可解性）；
  \item \textbf{复述}经典视觉 SLAM 五模块框架（传感器、前端 VO、后端、回环、建图）各自的输入、输出、职责与彼此的数据流；
  \item \textbf{写出} SLAM 的离散时间运动方程 $\mathbf{x}_k=f(\mathbf{x}_{k-1},\mathbf{u}_k,\mathbf{w}_k)$ 与观测方程 $\mathbf{z}_{k,j}=h(\mathbf{y}_j,\mathbf{x}_k,\mathbf{v}_{k,j})$，并能为平面运动、二维激光观测写出具体参数化；
  \item \textbf{解释} SLAM 为何是一个\emph{状态估计}问题，并据“线性/非线性、高斯/非高斯”把求解方法分类为滤波（KF/EKF/PF）与优化（图优化）两条主线；
  \item \textbf{区分}按传感器（视觉/激光/视觉-惯性）与按方法（特征法/直接法、滤波/优化）的 SLAM 分类，并对应到 MonoSLAM、PTAM、ORB-SLAM、VINS-Mono、LIO-SAM 等代表系统；
  \item \textbf{搭建}最小的 C++/CMake 工程，理解“源文件—库—头文件—链接—可执行程序”的工程链条，跑通第一个 “Hello SLAM”；
  \item \textbf{使用}本章的全书路线图，知道每一个后续问题（位姿如何表达与优化、相机如何成像、方程如何求解、前端/后端/回环/建图如何实现）将在哪一章解决。
\end{enumerate}

\subsection*{本章知识导航}
本章是一条“\emph{从一台想自己走路的机器人，到一套完整的 SLAM 知识地图}”的主线。先用直觉立起定位与建图的相互依赖，再把它拆成工程上的五个模块，接着把模块抽象成两条数学方程并归结为状态估计问题，最后落到分类、代表系统与编程环境。结构如下：

\begin{center}
\small
\begin{tikzpicture}[
  node distance=4mm and 7mm, >=Stealth,
  blk/.style={draw, rounded corners, align=center, font=\small, inner sep=3pt, fill=main!6, draw=main!60!black},
  grp/.style={draw, rounded corners, align=center, font=\small, inner sep=3pt, fill=second!6, draw=second!60!black},
  ln/.style={->, thick, gray!60!black}]
\node[blk] (why) {SLAM 是什么\\为什么需要};
\node[blk, right=of why] (frame) {五模块框架\\（工程图景）};
\node[blk, right=of frame] (math) {运动/观测方程\\（数学表述）};
\node[grp, below=8mm of why] (est) {状态估计\\滤波 vs 优化};
\node[grp, right=of est] (cls) {分类\\传感器/方法};
\node[grp, right=of cls] (sys) {代表系统\\ORB/VINS/LIO};
\node[blk, below=8mm of est, fill=third!8, draw=third!60!black] (code) {Hello SLAM\\CMake/库/链接};
\node[blk, right=of code, fill=third!8, draw=third!60!black] (map) {全书路线图\\（每章地图）};
\draw[ln] (why)--(frame); \draw[ln] (frame)--(math);
\draw[ln] (math.south) -- ++(0,-4.5mm) -| (est.north);
\draw[ln] (est)--(cls); \draw[ln] (cls)--(sys);
\draw[ln] (est.south)-- (code.north);
\draw[ln] (code)--(map);
\end{tikzpicture}
\end{center}

\noindent\textbf{推荐路径}：\cref{sec:intro-what,sec:intro-frame} 是任何读者的主干（建立“为什么需要 SLAM”与“五模块”两张图）；\cref{sec:intro-math,sec:intro-est} 是把工程图景数学化、并引出全书方法论的枢纽，后续状态估计、优化、滤波各章会反复回看；\cref{sec:intro-class,sec:intro-systems} 给出 SLAM 的分类坐标系与现代系统坐标，可作为速览；\cref{sec:intro-code} 是上机准备，建议照着敲一遍；\cref{sec:intro-roadmap} 是本章对全书的承诺——它把后面每一章都点到名。

\subsection*{前置知识桥接}
本章是全书的\emph{开篇}，不假设任何 SLAM 背景，只假设读者具备基本的线性代数、微积分与概率，以及一点 Linux/C++ 常识。它\emph{回避}了所有“细节怎么做”的问题——位姿如何表达与优化、相机如何把空间点投成像素、非线性最小二乘如何求解、特征如何提取匹配——并把这些问题逐一\emph{命名}并指向后续章节。读完本章，你将拥有一张全书地图，知道每一块拼图归属何处。

\begin{note}
\textbf{本章记号与约定.}\;
本章是全书提要，记号大多停留在抽象层面（用一般函数 $f,h$），与全书统一约定差异极小，登记如下。
离散时刻统一用 $k=1,\dots,K$（部分文献写 $t$，本书统一 $k$）；机器人位姿/状态记 $\mathbf{x}_k$（三维时 $\mathbf{x}_k\in\SE(3)$ 或其参数化，见 \cref{ch:rigid_body,ch:lie}）；路标/地图点记 $\mathbf{y}_j$（$j=1,\dots,N$；部分文献用 $\mathbf{m}_i$，本书沿用 $\mathbf{y}_j$，地图整体记 $\mathbf{m}$）；运动输入/控制记 $\mathbf{u}_k$；观测记 $\mathbf{z}_{k,j}$（$k$ 时刻对 $\mathbf{y}_j$ 的观测）。
运动噪声 $\mathbf{w}_k\sim\mathcal{N}(\mathbf{0},\boldsymbol{\Sigma}_w)$、观测噪声 $\mathbf{v}_{k,j}\sim\mathcal{N}(\mathbf{0},\boldsymbol{\Sigma}_v)$（文献处理不一：部分教材在引入方程时只写抽象噪声项 $\mathbf{w}_k,\mathbf{v}_{k,j}$、暂不指定其分布\cite{gaoxiang2019slam14}，而严谨的状态估计文献多在一开始即设零均值高斯\cite{barfoot2024state,thrun2005probabilistic}；本书统一在 \cref{eq:intro-motion,eq:intro-obs} 处即取零均值高斯，并在状态估计章兑现）；信息矩阵 $\boldsymbol{\Omega}=\boldsymbol{\Sigma}^{-1}$。观测索引集合 $\mathcal{O}$ 记录“哪个时刻看到了哪个路标”。本章不涉及旋转的具体表示与四元数，留到 \cref{ch:rigid_body,ch:lie}。
\end{note}

% ════════════════════════════════════════════════════════════════
\section{SLAM 是什么，为什么需要它}
\label{sec:intro-what}

\paragraph{动机：让机器人自己走路。}
设想我们组装一台叫“小萝卜”的机器人\cite{gaoxiang2019slam14}：身上有相机、轮子，还有一台塞进后备箱、方便随时调试的笔记本电脑。我们希望它具有\emph{自主运动能力}——这是扫地、搬东西等一切高级功能的前提。要移动就得有轮子和电机；但若不加规划与控制，机器人不知道自己要去哪里，会四处乱走，甚至撞墙损毁。要规划与控制，就得\emph{先感知周边环境}，于是我们在它脑袋上装了相机。这其实是在模仿人类：一个有眼睛、大脑和四肢的人，能在任意环境里轻松行走、探索。

\paragraph{两个最基本的问题。}
那么，为探索一个房间，机器人至少需要知道哪些事？答案是两件，且只有两件\cite{gaoxiang2019slam14}：

\begin{enumerate}
  \item \textbf{我在什么地方？}——这是\emph{定位}（Localization）。
  \item \textbf{周围环境是什么样？}——这是\emph{建图}（Mapping）。
\end{enumerate}

\noindent 定位与建图，可以看成感知的“内外之分”：作为一个“内外兼修”的机器人，一方面要明白\emph{自身的状态}（即位置），另一方面也要了解\emph{外在的环境}（即地图）。下面我们会看到，这两件事并不能各自独立地解决——它们彼此咬合，构成 SLAM 的核心矛盾。

\begin{definition}[同时定位与建图 SLAM]\label{def:intro-slam}
\emph{同时定位与建图}（Simultaneous Localization and Mapping, SLAM）指：一个携带传感器的机器人，在\textbf{没有任何关于自身位置的先验知识}的前提下，于运动过程中\emph{同时}估计自身的状态（轨迹）并构建所处环境的模型（地图）\cite{cadena2016past,durrantwhyte2006slam}。在最简单的情形，“状态”即机器人的位姿（位置与朝向），也可包含速度、传感器偏置、标定参数等；“地图”是对环境中所关心方面（路标、障碍物、表面）的表示。当传感器主要是相机时，称\emph{视觉 SLAM}；主要是激光时，称\emph{激光 SLAM}。
\end{definition}

\noindent 这一权威定义可与三本经典互相印证。Cadena 等\cite{cadena2016past} 写道：SLAM 是“对装有车载传感器的机器人状态的同时估计，以及对这些传感器所感知环境的模型（地图）的构建”。Durrant-Whyte 与 Bailey\cite{durrantwhyte2006slam} 强调它的“在线、无先验”特性：平台的整条轨迹与所有路标的位置都是\emph{在线估计}的，\emph{无须任何关于位置的先验知识}。SLAM Handbook 的视觉 SLAM 章\cite{carlone2026handbook} 则补上“实时”：在未知环境中实时地同时自定位与建图。把这三句话叠在一起，正是 \cref{def:intro-slam} 的全部内涵：未知、在线、同时、实时。

\subsection{为什么不能一劳永逸地“装个传感器”？}
\label{sec:intro-sensors}

\paragraph{解决定位的两条路。}
其实，让机器人知道“我在哪”有很多现成办法：地板铺导引线、墙上贴二维码、桌上放无线电定位器（很多仓储物流机器人就这么干）、室外装 GPS 接收器。按传感器“装在哪”，可把它们分成两大类\cite{gaoxiang2019slam14}，见 \cref{tab:intro-sensors}。

\begin{table}[H]\centering\small
\caption{两类传感器的对比：为什么 SLAM 偏要用“携带式”}
\label{tab:intro-sensors}
\begin{tabular}{p{0.16\textwidth} p{0.30\textwidth} p{0.18\textwidth} p{0.26\textwidth}}
\toprule
类别 & 典型设备 & 优点 & 局限 \\
\midrule
\textbf{安装于环境中}（外部） & 导轨/导引线、二维码标志、（室内）无线电定位、GPS 接收器 & 通常能\emph{直接测量}机器人位置，简单可靠 & \textbf{要求环境必须人工布置}，限制使用范围（室内无 GPS、园区难铺导轨） \\
\textbf{携带于机器人本体}（内部） & 轮式编码器、相机、激光雷达、IMU（惯性测量单元） & \textbf{不对环境提任何要求}，适用于未知环境 & 测量的是\emph{间接物理量}（编码器测轮转角、IMU 测角速度/加速度、相机/激光读外部观测），只能\emph{间接推算}位置 \\
\bottomrule
\end{tabular}
\end{table}

\paragraph{命门在“未知环境”四个字。}
环境类传感器虽然简单可靠，却约束了外部环境——一旦约束无法满足（贴不了二维码、铺不了导轨、室内收不到 GPS），定位就无从谈起，因此它无法提供\emph{普遍、通用}的解决方案。而携带式传感器“没有对环境提出任何要求”，使定位方案可适用于\emph{未知}环境\cite{gaoxiang2019slam14}。这正是 SLAM 的立足点：我们在 SLAM 中非常强调未知环境，所以理论上不限制机器人的使用环境，也就\emph{不能}假设像 GPS、导轨这样的外部设施都能顺利工作。于是，用携带式传感器完成定位与建图，成了我们重点关心的问题。

\begin{insight}
SLAM 的全部理论张力，都来自“未知环境 + 携带式传感器”这一组合。携带式传感器测的是\emph{间接量}（轮子转了多少、转了多快、看到了什么），机器人必须\emph{从这些间接量反推自己的位置和环境}。没有外部地标可依，它只能\emph{一边走、一边把环境建起来当作参照}——这就逼出了下面要讲的“鸡生蛋”循环，也逼出了整本书。
\end{insight}

\subsection{定位与建图的相互依赖：“鸡生蛋”问题}
\label{sec:intro-chicken-egg}

\paragraph{两个退化的子问题。}
为什么不把定位和建图\emph{分开}做？因为各自单独做，都要求对方已经做好。把它们写成概率语言看得最清楚\cite{durrantwhyte2006slam}。记 $\mathbf{x}_k$ 为车辆状态、$\mathbf{m}$ 为地图、$Z_{0:k}$ 与 $U_{0:k}$ 分别为到 $k$ 时刻为止的全部观测与控制，则：

\begin{itemize}
  \item \textbf{纯建图问题} = 计算 $P(\mathbf{m}\mid X_{0:k},Z_{0:k},U_{0:k})$，它\emph{假设车辆在所有时刻的位置 $X_{0:k}$ 已知（确定）}，仅由不同位置上的观测融合出地图；
  \item \textbf{纯定位问题} = 计算 $P(\mathbf{x}_k\mid Z_{0:k},U_{0:k},\mathbf{m})$，它\emph{假设地图 $\mathbf{m}$ 确定已知}，目标是估计车辆相对这些已知路标的位置。
\end{itemize}

\noindent 两个子问题各自把对方当成“已知”。可是在 SLAM 里，位姿与地图\emph{都未知}：要定位，需要地图作参照；要建图，又需要已知位姿才能把观测放到正确的位置上。这就是著名的“鸡生蛋”（chicken-and-egg）困境——也正是 SLAM 这个缩写里“Simultaneous（同时）”一词的由来：二者必须\emph{联合}估计。

\begin{pitfall}[概念误区：联合后验可以“拆开”分别算]
新手最自然的念头是：既然要同时估计 $\mathbf{x}_k$ 和 $\mathbf{m}$，那就分别算 $P(\mathbf{x}_k\mid \mathbf{z}_k)$ 与 $P(\mathbf{m}\mid \mathbf{z}_k)$ 再拼起来，即假设
\begin{equation}
  P(\mathbf{x}_k,\mathbf{m}\mid \mathbf{z}_k)\;\overset{?}{=}\;P(\mathbf{x}_k\mid \mathbf{z}_k)\,P(\mathbf{m}\mid \mathbf{z}_k).
  \label{eq:intro-bad-partition}
\end{equation}
这是\textbf{错的}。观测模型 $P(\mathbf{z}_k\mid \mathbf{x}_k,\mathbf{m})$ \emph{同时}依赖车辆与路标位置，联合后验不能如此分解；早期关于一致性建图的研究就已指出，\textbf{这样的分割会得到不一致的估计}\cite{durrantwhyte2006slam}。换言之，位姿误差与地图误差是\emph{纠缠}在一起的，必须一起算。
\end{pitfall}

\paragraph{为什么纠缠，又为什么可解。}
误差为何纠缠？因为路标位置的估计误差有一个\emph{共同来源}——“观测这些路标时，机器人到底在哪儿”的认知误差\cite{durrantwhyte2006slam}。同一时刻看到的若干路标，其位置估计都受同一个车辆位置误差的影响，于是它们\emph{高度相关}：任意两个路标的\emph{相对}位置 $\mathbf{m}_i-\mathbf{m}_j$ 可被高精度获知，即使单个路标的\emph{绝对}位置 $\mathbf{m}_i$ 仍很不确定（概率上：联合密度 $P(\mathbf{m}_i,\mathbf{m}_j)$ 很尖锐，而边缘密度 $P(\mathbf{m}_i)$ 很弥散）。

这层相关性恰恰是 SLAM \emph{可解}的关键。Durrant-Whyte 与 Bailey\cite{durrantwhyte2006slam} 称之为 SLAM 最重要的洞见：

\begin{theorem}[SLAM 的单调收敛性（线性高斯情形）]\label{thm:intro-convergence}
随着观测不断增多，各路标位置估计之间的\emph{相关性单调增加}；等价地，所有路标的联合密度 $P(\mathbf{m})$ 随观测增多而\emph{单调变尖锐}，路标相对位置的认知\emph{只会改善、永不发散}，且与机器人如何运动无关。在 EKF-SLAM 框架下，这表现为地图协方差矩阵行列式\emph{单调收敛}到一个由\emph{初始}不确定性（机器人初始位置与观测精度）决定的下界。该结论已在\emph{线性高斯}情形被严格证明；一般概率情形的正式证明仍是开放问题\cite{durrantwhyte2006slam}。
\end{theorem}

\begin{insight}
\textbf{弹簧网络类比}\cite{durrantwhyte2006slam}：把所有路标想成由弹簧连接的网络，弹簧的“刚度”就是路标间的相关性。机器人每观测一次，相当于把若干弹簧拉紧一点；它在环境里来回穿行，弹簧（单调地）越来越硬，极限下整张网变成一块\emph{刚性的相对地图}。此时机器人相对这张地图的定位精度，只受地图本身与相对测量传感器的质量所限。这幅图景告诉我们：鸡生蛋的循环不是死结——\emph{每一次联合更新都同时让位姿和地图变好一点}，循环往复，终归收敛。
\end{insight}

\subsection{什么时候\texorpdfstring{\emph{不}}{不}需要 SLAM？为什么仍然需要它？}
\label{sec:intro-need}

\paragraph{反面：地图已知时不必 SLAM。}
若一组路标的位置\emph{先验已知}（例如工厂地面人工布设了信标地图，或机器人有可靠 GPS——可把 GPS 卫星看作位置已知的移动信标），且机器人能可靠地相对这些已知路标定位，那么它\emph{不需要} SLAM——这退化为纯定位问题\cite{cadena2016past}。SLAM 的流行恰与\emph{室内移动机器人}应用的兴起相伴：室内排除了 GPS，SLAM 给出了一种\emph{无须人工布设定位基础设施}即可运行的方案。

\paragraph{那自主机器人“真的”需要 SLAM 吗？}
这是一个被反复追问的问题，Cadena 等\cite{cadena2016past} 给出三重回答，其关键词是\emph{回环}（loop closure）：

\begin{quote}
“SLAM 旨在构建环境的\emph{全局一致}表示，它同时利用自运动测量与回环。这里的关键词是‘回环’：\textbf{如果我们放弃回环，SLAM 就退化为里程计。}”\cite{cadena2016past}
\end{quote}

\begin{enumerate}
  \item \textbf{SLAM 研究本身产出了当前最先进的里程计}：现代视觉-惯性里程计的漂移已小于轨迹长度的 $0.5\%$；从这个意义上说，\emph{视觉-惯性导航（VIN）就是禁用了回环/地点识别模块的简化 SLAM}\cite{cadena2016past}。
  \item \textbf{回环揭示环境的真实拓扑}：只做里程计而忽略回环，机器人会把世界理解成一条“无限长的走廊”，从起点不断探索“新”区域；而一次回环事件告诉它这条走廊其实\emph{自相交}，从而发现地图中的“捷径”。为什么不干脆丢掉度量信息、只做地点识别？因为\emph{度量信息让地点识别更简单、更鲁棒}——度量重建能预测哪里有回环机会，并丢弃虚假回环。SLAM 因此是抵御\emph{错误数据关联}与\emph{感知混叠}（perceptual aliasing，外观相似实为不同地点）的天然防线\cite{cadena2016past}。
  \item \textbf{许多应用显式或隐式要求全局一致地图}：军事/民用的探索与汇报、桥梁与建筑的结构巡检，都需要一张能拼到一起、不自相矛盾的地图\cite{cadena2016past}。
\end{enumerate}

\paragraph{那 SLAM “解决了”吗？}
这个问题只有针对具体的\emph{机器人/环境/性能}三元组才有意义\cite{cadena2016past}：要指定机器人（运动类型、传感器、算力）、环境（平面或三维、自然或人工路标、动态程度、对称性与感知混叠风险）、性能（精度、地图类型、成功率、延迟、运行时长、地图规模）。例如：二维室内 + 轮式编码器 + 激光、精度优于 $10\,$cm、低失败率，已可认为\emph{基本解决}（工业界如 Kuka 导航方案）；慢速视觉 SLAM（火星车、家用机器人）与视觉-惯性里程计是\emph{成熟的研究领域}；而高速、高度动态、需高频闭环控制的场景，仍需大量基础研究。本章末的现代系统坐标（\cref{sec:intro-systems}）正是为这张“成熟度地图”提供锚点。

\begin{practice}
\begin{enumerate}
  \item 用自己的话说明：为什么“先建好图、再用图定位”和“先定好位、再用位姿建图”都行不通？把它们对应到 \cref{sec:intro-chicken-egg} 的两个退化子问题。
  \item 举出两个\emph{不需要} SLAM 的真实场景，并说明它们各自靠什么提供了“地图已知”或“位置已知”。
  \item 解释“放弃回环，SLAM 就退化为里程计”这句话：里程计相比完整 SLAM 缺了什么能力？
\end{enumerate}
\end{practice}

\subsection{SLAM 简史：三个时代}
\label{sec:intro-history}

\paragraph{起源：1986 与 1995。}
概率 SLAM 问题的起源公认是 \emph{1986 年旧金山 IEEE 机器人与自动化会议}（ICRA）：当时概率方法刚被引入机器人与人工智能，Peter Cheeseman、Jim Crowley、Hugh Durrant-Whyte 等人在会上讨论如何一致地建图\cite{durrantwhyte2006slam}。随后几年的关键论文（Smith \& Cheeseman 1986、Durrant-Whyte 1988）建立了描述路标间关系、操纵几何不确定性的\emph{统计基础}，并指出不同路标的位置估计必然\emph{高度相关}、相关性随观测增长——这正是 \cref{thm:intro-convergence} 的雏形。Smith、Self \& Cheeseman 进一步证明：移动机器人取相对路标观测时，所有路标估计因\emph{车辆位置的共同误差}而必然相互相关，故一致的完整解需要一个由\emph{车辆位姿 + 每个路标位置}组成的联合状态向量，每次观测后整体更新——而其计算量随路标数\emph{平方}增长\cite{durrantwhyte2006slam}。\textbf{“SLAM” 这个缩写本身，连同收敛性结论，首次系统提出于 1995 年的国际机器人研究研讨会（ISRR）}\cite{durrantwhyte2006slam}（此前也称 CML，concurrent mapping and localization）。

\paragraph{三个时代。}
Cadena 等\cite{cadena2016past} 把此后三十年划为三个时代，见 \cref{tab:intro-eras}。

\begin{table}[H]\centering\small
\caption{SLAM 的三个时代（据 Cadena 等\cite{cadena2016past}）}
\label{tab:intro-eras}
\begin{tabular}{p{0.20\textwidth} p{0.12\textwidth} p{0.56\textwidth}}
\toprule
时代 & 年代 & 主要特征 \\
\midrule
\textbf{古典时代}\newline(Classical age) & 1986--2004 & 确立 SLAM 的主要\emph{概率表述}：扩展卡尔曼滤波、Rao-Blackwell 化粒子滤波、最大似然估计；勾勒了\emph{效率}与\emph{鲁棒数据关联}两大基本挑战。 \\
\textbf{算法分析时代}\newline(Algorithmic-analysis age) & 2004--2015 & 研究 SLAM 的\emph{基本性质}：可观性、收敛性、一致性；认识到\emph{稀疏性}对高效求解的关键作用；主要的\emph{开源 SLAM 库}在此期间出现。 \\
\textbf{鲁棒感知时代}\newline(Robust-perception age) & 2015--至今 & 追求：(1) \emph{鲁棒性}（低失败率、长时段、广环境、自检自调，对抗“手工调参的诅咒”）；(2) \emph{高层理解}（超越基础几何，理解语义、物理、可供性）；(3) \emph{资源感知}；(4) \emph{任务驱动的感知}。 \\
\bottomrule
\end{tabular}
\end{table}

\noindent 这条时间线不是掌故，而是全书的隐藏脉络：古典时代的滤波与概率表述对应 \cref{ch:eskf} 的滤波器与 \cref{ch:slam_est} 的状态估计；算法分析时代的“稀疏性”是 \cref{ch:nlopt} 后端优化高效求解的命根；鲁棒感知时代的诉求，则是 \cref{ch:loop} 回环鲁棒化、\cref{ch:vio,ch:lidar_slam} 多传感器融合的现实驱动。

% ════════════════════════════════════════════════════════════════
\section{看清相机：单目、双目与 RGB-D}
\label{sec:intro-camera}

\paragraph{动机：相机就是视觉 SLAM 的“眼睛”。}
既然视觉 SLAM 用相机解决定位与建图（\cref{def:intro-slam}），那就得先弄清相机给我们的是什么数据。SLAM 用的相机不同于单反：它更简单、通常不带昂贵镜头，以一定速率（普通约 $30$ 帧/秒）连续拍摄环境，形成\emph{连续的视频流}\cite{gaoxiang2019slam14}。按工作方式，主流相机分单目、双目、深度三类（此外还有全景相机、事件相机等尚未成主流的特殊种类）。它们的差别，全在于\emph{能不能直接拿到深度}。

\subsection{单目相机：从一张照片到尺度不确定性}
\label{sec:intro-mono}
只用一个摄像头做 SLAM，称\emph{单目 SLAM}（Monocular SLAM）。它结构特别简单、成本特别低，数据就是一张张照片。但照片的本质，是某个场景在相机成像平面上留下的一个\emph{投影}——\textbf{以二维记录三维，丢掉了一个维度：深度（距离）}\cite{gaoxiang2019slam14}。因此，\emph{无法从单张图片算出物体与相机的距离}：手掌上的小人是真人还是模型？仅凭一张图无法判断。换言之，单张图里无法确定物体的真实大小——“很大但很远”与“很近但很小”，因近大远小的透视关系，在图像中可能等大。

要恢复三维，就必须\emph{改变视角、移动相机}。相机一移动，就能估计它的\emph{运动}（Motion），同时估计场景物体的远近大小即\emph{结构}（Structure）。机制是\emph{视差}（Disparity）：相机移动时，近处物体在图像上移动快、远处慢、极远处（如日月）看上去不动；由视差就能\emph{定量}判断远近。但这里埋着单目相机的一个根本性麻烦：

\begin{theorem}[单目的尺度不确定性]\label{thm:intro-scale}
即便恢复了物体间的相对远近，那也只是\emph{相对}值。若把相机的运动与场景的大小\emph{同时}放大任意倍，单目相机看到的图像序列\emph{完全一样}。因此单目 SLAM 估计的轨迹与地图，与真实世界相差一个未知的整体比例因子，称\emph{尺度}（Scale）；\textbf{仅凭单目图像无法确定真实尺度}，称\emph{尺度不确定性}（Scale Ambiguity）\cite{gaoxiang2019slam14}。
\end{theorem}

\noindent 于是单目 SLAM 有两大根本麻烦：\emph{必须平移之后才能算深度}，以及\emph{无法确定真实尺度}；其共同根源是“通过单张图像无法确定深度”。为绕开它，人们引入了能直接给出深度的双目相机与深度相机。

\subsection{双目相机与深度相机：直接拿到深度的两种办法}
\label{sec:intro-stereo-rgbd}
\paragraph{双目相机（Stereo）。}
双目相机由两个单目相机组成，两者间距称\emph{基线}（Baseline，已知）。通过基线，可估计每个像素的空间位置——正如人眼靠左右两图的差异判断远近\cite{gaoxiang2019slam14}。它可拓展为多目相机（本质相同），\emph{室内外皆可用}（不依赖其他传感设备，仅比较左右图）。代价是：计算机做双目需\emph{大量计算}才能（不太可靠地）估计每个像素的深度；可测深度范围\emph{与基线相关}——基线越大、能测越远（无人车上的双目通常很大）；视差计算极耗算力，往往需 GPU/FPGA 加速才能实时输出整张图的距离。\textbf{计算量，是双目的主要问题之一}。

\paragraph{深度相机（RGB-D）。}
深度相机于 2010 年前后兴起，最大特点是通过\emph{红外结构光}或\emph{飞行时间}（Time-of-Flight, ToF）原理，像激光那样\emph{主动发射光并接收返回光}，\textbf{物理地}测出距离——相比双目的软件计算，可节省大量算力\cite{gaoxiang2019slam14}。常见型号有 Kinect / Kinect V2、Xtion Pro Live、RealSense 等（部分手机也用它做人脸识别）。但它存在测量范围窄、噪声大、视野小、易受日光干扰、无法测量透射材质等问题，因此在 SLAM 中\emph{主要用于室内，室外较难应用}。RGB-D 的成像与深度模型将在 \cref{ch:camera} 详述。

\begin{table}[H]\centering\small
\caption{三类相机对比}
\label{tab:intro-cameras}
\begin{tabular}{p{0.12\textwidth} p{0.10\textwidth} p{0.16\textwidth} p{0.46\textwidth}}
\toprule
类型 & 摄像头数 & 能否直接测深度 & 原理 / 特点 \\
\midrule
单目 Monocular & 1 & 否 & 二维记录三维、丢失深度；需平移才能算深度；存在尺度不确定性 \\
双目 Stereo & 2（或多目） & 能（软件计算） & 已知基线、比较左右图差异；室内外皆可；计算量大、需 GPU/FPGA \\
深度 RGB-D & 多个（含主动发射） & 能（物理测量） & 红外结构光或 ToF 主动测距；省算力；主要室内 \\
\bottomrule
\end{tabular}
\end{table}

\noindent 想象相机在场景中运动，得到一系列连续变化的图像。视觉 SLAM 的目标，就是\emph{通过这样一些图像，进行定位与地图构建}\cite{gaoxiang2019slam14}。但这不是某个“输入图像就吐出定位与地图”的单一算法，而需要\emph{一整套完善的算法框架}。经研究者长期努力，这套框架已经相当成熟——下一节我们就来看它。

% ════════════════════════════════════════════════════════════════
\section{经典视觉 SLAM 框架：五个模块}
\label{sec:intro-frame}

\paragraph{动机：把一个大问题拆成可分工的零件。}
SLAM 是个大问题，一口吃不下。过去十几年的研究，把视觉 SLAM 沉淀成了一套\emph{模块化}框架\cite{gaoxiang2019slam14}：图像从传感器流入，经过几个各司其职的模块，最终流出一条全局一致的轨迹和一张地图。理解这五个模块的\emph{职责}与彼此的\emph{数据流}，就理解了视觉 SLAM 的工程骨架——它同时也是本书后半部分的目录。

\subsection{五模块总览与数据流}
\label{sec:intro-pipeline}
整个视觉 SLAM 流程包含以下步骤\cite{gaoxiang2019slam14}：

\begin{enumerate}
  \item \textbf{传感器信息读取}（Sensor data）。视觉 SLAM 中主要是相机图像的读取与预处理；机器人上还可能有码盘、惯性传感器等信息的读取与\emph{同步}。
  \item \textbf{前端视觉里程计}（Visual Odometry, VO）。任务是估算\emph{相邻图像间相机的运动}，以及局部地图的样子。VO 又称\emph{前端}（Front End）。
  \item \textbf{后端（非线性）优化}（Optimization）。后端接收不同时刻 VO 测得的相机位姿，以及回环检测的信息，对它们做优化，得到\emph{全局一致}的轨迹与地图。因接在 VO 之后，又称\emph{后端}（Back End）。
  \item \textbf{回环检测}（Loop Closure Detection）。判断机器人\emph{是否到达过先前的位置}；若检测到回环，把信息提供给后端处理。
  \item \textbf{建图}（Mapping）。根据估计的轨迹，建立\emph{与任务要求对应}的地图。
\end{enumerate}

\noindent 这五个模块的连接关系如 \cref{fig:intro-pipeline}：传感器数据流入前端，前端把\emph{相机位姿初值 + 局部结构}交给后端，同时把图像流送给回环检测；回环检测把“A 与 B 是同一个点”的约束送给后端；后端融合前端位姿与回环约束做\emph{全局优化}，输出全局一致的轨迹与地图；建图据此生成最终地图。现代综述\cite{cadena2016past} 还强调：\emph{后端会反馈信息给前端}，用于回环的检测与验证（见 \cref{sec:intro-frontback}）。

\begin{figure}[ht]
\centering
\begin{tikzpicture}[
  >=Stealth, node distance=10mm,
  mod/.style={draw, rounded corners, align=center, font=\small, inner sep=4pt, fill=main!7, draw=main!60!black, minimum height=11mm, text width=22mm},
  lp/.style={draw, rounded corners, align=center, font=\small, inner sep=4pt, fill=second!8, draw=second!60!black, minimum height=11mm, text width=22mm},
  flow/.style={->, thick, gray!55!black},
  fb/.style={->, thick, second!60!black, dashed}]
\node[mod] (sen) {传感器数据\\Sensor};
\node[mod, right=of sen] (vo) {前端 VO\\Front End};
\node[mod, right=of vo] (bk) {后端优化\\Back End};
\node[mod, right=of bk] (mp) {建图\\Mapping};
\node[lp, below=12mm of vo] (loop) {回环检测\\Loop Closure};
\draw[flow] (sen) -- node[above, font=\scriptsize]{图像/帧} (vo);
\draw[flow] (vo) -- node[above, font=\scriptsize, align=center]{位姿初值\\+局部地图} (bk);
\draw[flow] (bk) -- node[above, font=\scriptsize, align=center]{全局一致\\轨迹+地图} (mp);
\draw[flow] (vo.south) -- node[left, font=\scriptsize]{图像} (loop.north);
\draw[flow] (loop.east) -| node[pos=0.25, below, font=\scriptsize]{回环约束(A=B)} (bk.south);
\draw[fb] (bk.north) .. controls +(0,7mm) and +(0,7mm) .. node[above, font=\scriptsize]{反馈：回环验证} (vo.north);
\end{tikzpicture}
\caption{经典视觉 SLAM 框架的五模块与数据流。实线为前向数据流，虚线为后端到前端的反馈（用于回环检测与验证）。仿 \cite{gaoxiang2019slam14,cadena2016past} 重绘。}
\label{fig:intro-pipeline}
\end{figure}

\begin{table}[H]\centering\small
\caption{五模块的输入、输出、职责与去向章节}
\label{tab:intro-modules}
\begin{tabular}{p{0.11\textwidth} p{0.15\textwidth} p{0.17\textwidth} p{0.20\textwidth} p{0.16\textwidth}}
\toprule
模块 & 输入 & 输出 & 核心职责（关键问题） & 详述章 \\
\midrule
传感器读取 & 相机图像（+码盘/IMU 原始数据） & 同步、预处理后的帧 & 读取与同步多源数据（时间同步） & \cref{ch:camera,ch:imu} \\
前端 VO & 相邻（或近邻 $5\!\sim\!10$）帧图像 & 相邻帧相机运动 + 局部结构 & 估相邻帧间运动、恢复局部结构（\emph{累积漂移}） & \cref{ch:vo} \\
后端优化 & 各时刻 VO 位姿 + 回环约束 & 全局一致轨迹 + 地图 + 不确定性 & 从带噪数据估系统状态及其不确定性（MAP） & \cref{ch:slam_est,ch:nlopt,ch:eskf} \\
回环检测 & 图像（当前 vs 历史） & “A 与 B 同一点”的约束 & 判断是否回到旧地点、提供约束（图像相似性） & \cref{ch:loop} \\
建图 & 估计的轨迹（+观测） & 任务相应的地图 & 按应用构建地图（无固定形式） & \cref{ch:dense,ch:pointcloud} \\
\bottomrule
\end{tabular}
\end{table}

\noindent 需要强调框架的\emph{适用边界}：经典视觉 SLAM 框架是过去十几年的研究成果，但它的“成熟”有前提——若把工作环境限定在\emph{静态、刚体、光照变化不明显、没有人为干扰}的场景，这种场景下的 SLAM 技术已相当成熟\cite{gaoxiang2019slam14}。一旦环境动态、光照剧变或有人干扰，框架就会受挑战，这正是鲁棒感知时代（\cref{tab:intro-eras}）的攻关方向。

\subsection{前端视觉里程计：能走多远，又为何走偏}
\label{sec:intro-vo}
\paragraph{VO 在做什么。}
VO 关心\emph{相邻图像之间的相机运动}，最简单就是两张图像之间的运动关系：人一眼能看出右图是左图“向左转了一点”，却\emph{说不清究竟转了多少度、平移了多少厘米}——直觉对具体数字不敏感，而计算机必须\emph{精确测量}\cite{gaoxiang2019slam14}。难点在于：图像在计算机里只是一个\emph{数值矩阵}，矩阵里表达了什么，计算机毫无概念；视觉 SLAM 中它只能看到一个个像素，知道它们是某些空间点在成像平面上投影的结果。因此，\textbf{为定量估计相机运动，必须先了解相机与空间点的几何关系}（这正是 \cref{ch:camera} 相机模型、\cref{ch:vo} VO 实现要解决的）。

它之所以叫“里程计”，是因为像汽车里程计一样\emph{只计算相邻时刻的运动、与过去信息无关}（像只有短时记忆的物种；不限于两帧，可取 $5\!\sim\!10$ 帧）。它的输入是相邻帧图像，输出是相邻时刻的相机运动（位姿增量）加局部场景结构。有了 VO，似乎 SLAM 就解决了一半：把相邻时刻的运动一段段“串”起来，就得到轨迹（定位）；据每时刻相机位置算出各像素对应空间点的位置，就得到地图。

\paragraph{反面：只靠 VO 为什么不行——累积漂移。}
可是只靠 VO 串轨迹，会出大问题。VO 每次估计都带误差，\emph{先前时刻的误差会传递到下一时刻}，经过一段时间，轨迹就不再准确——这叫\emph{累积漂移}（Accumulating Drift）\cite{gaoxiang2019slam14}。

\begin{example}[累积漂移的一度之差]\label{ex:intro-drift}
设机器人先向左转 $90^\circ$，再向右转 $90^\circ$，理应回到原朝向。但若 VO 把第一个 $90^\circ$ 估成了 $89^\circ$，那么右转之后，估计的朝向就回不到原点；更糟的是，\emph{即便此后每一步都估计得完全准确，整条轨迹也都会带上这 $-1^\circ$ 的误差}。后果是无法建立一致的地图：原本笔直的走廊会被估成斜的，原本 $90^\circ$ 的直角不再是直角。
\end{example}

\noindent 这正是“放弃回环，SLAM 就退化为里程计”（\cref{sec:intro-need}）的另一面：纯里程计无法自我纠错。要消除漂移，需要两种技术联手——\emph{后端优化}与\emph{回环检测}：回环检测负责把“机器人回到了原始位置”这件事检测出来，后端优化据此\emph{校正整条轨迹的形状}。下面两节正讲它们。

\subsection{后端优化：与噪声共舞的状态估计}
\label{sec:intro-backend}
\paragraph{后端到底优化什么。}
后端优化主要处理 SLAM 过程中的\emph{噪声}问题\cite{gaoxiang2019slam14}：再精确的传感器也带噪声（便宜的误差大、贵的小、有的还受磁场温度影响）。除了“如何从图像估计相机运动”，我们还得关心：这个估计\emph{带多大噪声}、噪声\emph{如何从上一时刻传到下一时刻}、我们对当前估计\emph{有多大自信}。后端要回答的，正是：如何从这些带噪数据中估计整个系统的\emph{状态}及该估计的\emph{不确定性}——这称为\emph{最大后验概率估计}（Maximum-a-Posteriori, MAP）。这里的状态\emph{既包括机器人轨迹，也包括地图}。

\paragraph{前端与后端的分工。}
前端给后端提供\emph{待优化的数据及其初始值}；后端负责整体优化，\emph{往往只面对数据，不必关心数据来自什么传感器}\cite{gaoxiang2019slam14}。在视觉 SLAM 中，前端更接近计算机视觉（图像特征提取与匹配），后端主要是\emph{滤波与非线性优化}算法。这条分工线极其重要：它让后端成为一个通用的“状态估计引擎”，无论喂进来的是相机、激光还是 IMU 的约束，都能统一处理——这也是 \cref{ch:slam_est,ch:nlopt} 把后端抽象成纯数学问题来讲的根据。

\paragraph{历史的回响：SLAM 本就是状态估计。}
今天我们称为“后端优化”的部分，在很长一段时间里就\emph{直接}被称为“SLAM 研究”\cite{gaoxiang2019slam14}。早期 SLAM 本质上是一个状态估计问题；最早的一系列论文把它称作“空间状态不确定性的估计”（Spatial Uncertainty），这个名字精准反映了 SLAM 的本质——\emph{对运动主体自身与周围环境的空间不确定性的估计}。为此我们需要状态估计理论，把定位与建图的不确定性表达出来，再用\emph{滤波器或非线性优化}估计状态的均值与不确定性（方差）。这两条路线，正是 \cref{sec:intro-est} 要展开、并贯穿 \cref{ch:slam_est,ch:nlopt,ch:eskf} 的方法论主线。

\subsection{回环检测：把漂移“拉”回来}
\label{sec:intro-loop}
\paragraph{为什么需要回环。}
回环检测（又称闭环检测）主要解决\emph{位置估计随时间漂移}的问题\cite{gaoxiang2019slam14}。思路很朴素：机器人运动一段时间后\emph{实际}回到了原点，但因漂移，它\emph{估计}的位置没回到原点；若有办法让它意识到“我回到了原点”，再把位置估计“拉”过去，漂移就被消除了。回环检测与定位、建图都密切相关——\emph{地图存在的主要意义，就是让机器人知晓自己到过的地方}；因此我们需要让机器人具备\emph{识别到过场景}的能力。

\paragraph{怎么做回环。}
可以在环境里设标志物（如二维码），看到即知回到原点——但那又成了\emph{环境中的传感器}，对应用环境作了限制（不能贴二维码怎么办）。我们更希望用\emph{携带}的传感器，即\emph{图像本身}：通过判断\emph{图像间的相似性}完成回环检测（正如人看到两张相似的图片，很容易认出它们来自同一地方）。于是：\textbf{视觉回环检测实质上是一种计算图像数据相似性的算法}\cite{gaoxiang2019slam14}；图像信息非常丰富，这让正确检测回环的难度降低了不少。检测到回环后，把“A 与 B 是同一个点”的信息告诉后端，后端据此把轨迹和地图调整到符合回环的样子。\emph{充分而正确}的回环检测，能消除累积误差、得到全局一致的轨迹与地图。这套相似性计算（词袋模型等）将在 \cref{ch:loop} 详讲。

\subsection{建图：没有标准答案的模块}
\label{sec:intro-mapping}
\paragraph{地图随任务而定。}
建图就是构建地图的过程。但地图是对环境的\emph{描述}，而描述\emph{并不固定，要视 SLAM 应用而定}\cite{gaoxiang2019slam14}：家用扫地机器人在低矮平面运动，只需\emph{二维地图}标记可通过/障碍；相机有 $6$ 自由度运动，至少需\emph{三维地图}；有时想要带纹理三角面片的漂亮重建，有时只需知道“A 到 B 可通过、B 到 C 不行”，有时甚至不需要地图（或地图由他人提供，如行驶车辆用已绘的当地地图）。因此，相比 VO、后端、回环，\textbf{建图没有固定的形式与算法}——一组空间点、一个漂亮的 3D 模型、一张标着城市村庄铁路河道的图片，都可叫地图。大体可分\emph{度量地图}与\emph{拓扑地图}两类。

\paragraph{度量地图：稀疏与稠密。}
度量地图（Metric Map）强调\emph{精确表示物体的位置关系}，常用稀疏（Sparse）与稠密（Dense）划分，见 \cref{tab:intro-maps}。稀疏地图做了一定抽象，只选有代表意义的东西作为\emph{路标}（Landmark），由路标组成（非路标部分忽略）——\emph{定位用稀疏路标地图就够了}。稠密地图着重建模\emph{所有看到的东西}，按某分辨率由许多小块组成：二维是许多小格子（Grid），三维是许多小方块（体素，Voxel）；每个小块通常含\emph{占据、空闲、未知}三种状态，查询某空间位置即知能否通过——可用于 A*、D* 等导航算法。\emph{导航}往往需要稠密地图（否则会撞上两路标之间的墙），但代价是要存每个格点的状态、\emph{耗大量存储}（多数细节其实无用），且大规模时有\emph{一致性}问题（很小的转向误差就可能让两间屋子的墙重叠，使地图失效）。

\paragraph{拓扑地图：只关心连通性。}
拓扑地图（Topological Map）相比度量地图的精确性，更强调\emph{地图元素之间的关系}：它是一个\emph{图}（Graph），由节点和边组成，只考虑节点间的\emph{连通性}（只关注 A、B 是否连通，不关心如何从 A 到 B）。它放松了对精确位置的要求、去掉细节，是\emph{更紧凑}的表达。缺点是不擅长表达复杂结构的地图；如何把地图分割成节点与边、如何用拓扑地图做导航与路径规划，仍是有待研究的开放问题。建图的稠密重建、八叉树、TSDF 等将在 \cref{ch:dense,ch:pointcloud} 展开，而基于地图的导航与路径搜索（A*、D* 等）属规划范畴，见 \cref{ch:planning}。

\begin{table}[H]\centering\small
\caption{度量地图：稀疏 vs 稠密}
\label{tab:intro-maps}
\begin{tabular}{p{0.10\textwidth} p{0.32\textwidth} p{0.14\textwidth} p{0.30\textwidth}}
\toprule
子类 & 表示 & 适用 & 代价 / 问题 \\
\midrule
稀疏 Sparse & 由\emph{路标}组成，忽略非路标部分 & 定位 & ——（紧凑） \\
稠密 Dense & 二维=许多小格子 Grid；三维=许多体素 Voxel；每块含占据/空闲/未知 & 导航（A*、D*） & 耗大量存储；大规模时有一致性问题 \\
\bottomrule
\end{tabular}
\end{table}

\subsection{现代视角：前端/后端二分与并行流水线}
\label{sec:intro-frontback}
\paragraph{把五模块归并成两半。}
前一节把经典视觉 SLAM 框架按五模块划分（\cref{sec:intro-pipeline}），这种分法便于初学、与各模块的工程实现一一对应\cite{gaoxiang2019slam14}；现代综述则常统一为\emph{前端/后端二分}\cite{cadena2016past}：\textbf{前端}把原始传感器数据抽象成便于估计的模型，\textbf{后端}在这些抽象数据上做推断。VO 与回环检测属前端，优化属后端，建图依赖后端输出——两种视角等价。为什么需要前端？因为原始数据（如图像像素强度、单束激光）常难以直接写成状态的解析函数，故在后端之前先由前端完成三件事\cite{cadena2016past}：(1) \emph{特征提取}（视觉中提取少量可区分点的像素位置，使后端易于建模）；(2) \emph{数据关联}（把每个观测关联到一个路标子集，使观测可写成 $\mathbf{z}_k=h_k(\mathbf{x}_k)+\boldsymbol{\epsilon}_k$）；(3) 为后端的非线性优化提供\emph{初值}。

\begin{insight}
数据关联又分两种，理解它就理解了回环的位置\cite{cadena2016past}：\textbf{短期数据关联}（short-term）关联\emph{连续测量}中对应的特征（如跟踪相邻两帧里同一个 3D 点），这是 VO 在做的；\textbf{长期数据关联}（long-term）把\emph{新测量关联到很早以前的路标}，这正是\emph{回环}（loop closure）。换言之，VO 与回环检测做的是同一件事——数据关联——只是时间跨度不同。这也解释了为何后端常\emph{反馈}信息给前端：用当前的状态估计去帮助前端判断“这次回环是真是假”。
\end{insight}

\paragraph{并行流水线。}
SLAM Handbook 的视觉 SLAM 章\cite{carlone2026handbook} 把现代完整系统归纳为三个互补子功能：\emph{里程计前端 + 建图后端 + 回环与重定位}。前端 VO 有两条路线——\emph{特征法}（检测特征点、跨图像求对应、最小化重投影误差，末阶段类似经典光束法平差）与\emph{直接法}（直接对光度损失关于相机运动与结构优化）；后端用光束法平差或位姿图优化做全局优化；回环与重定位则在跟踪丢失或重访旧地时校正漂移。值得注意的是，tracking、mapping、optimization 等组件常\emph{并行运行}以提效（这一多线程思想源自 PTAM，并在 ORB-SLAM 中发扬光大，见 \cref{sec:intro-systems} 的现代系统）。这些前端实现的细节将分别在 \cref{ch:vo}（特征法/直接法）与 \cref{ch:loop}（回环）落地。

\begin{practice}
\begin{enumerate}
  \item 照 \cref{fig:intro-pipeline} 自己默画一遍五模块数据流，标出每条箭头上流动的是什么数据。
  \item \cref{ex:intro-drift} 中若第一个 $90^\circ$ 估成了 $91^\circ$，画出机器人“以为自己在哪”与“实际在哪”的差异；说明为何后续即使无误差，偏差也不会消失。
  \item 解释“VO 与回环检测都是数据关联，只是时间跨度不同”。哪一种对应短期、哪一种对应长期？
  \item 给三个任务（扫地机器人、桥梁结构巡检、室内送餐机器人），各推荐一种地图类型（稀疏/稠密、二维/三维、度量/拓扑），并说明理由。
\end{enumerate}
\end{practice}

% ════════════════════════════════════════════════════════════════
\section{SLAM 问题的数学表述}
\label{sec:intro-math}

\paragraph{动机：把工程图景翻译成可计算的方程。}
五模块框架给了我们工程图景，但要让计算机去解，必须把“机器人带着传感器在未知环境里运动”这句话\emph{翻译成数学}。这一节就完成这次翻译——得到两条基本方程，并由它们看清 SLAM 在数学上究竟是什么类型的问题。

\subsection{离散时间、状态量与路标量}
\label{sec:intro-discrete}
相机在某些时刻采集数据，故我们只关心这些时刻的位置与地图，把连续时间运动\emph{离散化}为时刻 $k=1,\dots,K$\cite{gaoxiang2019slam14}。用 $\mathbf{x}_k$ 表示机器人在 $k$ 时刻的位置，各时刻位置
\begin{equation}
  \mathbf{x}_1,\ \mathbf{x}_2,\ \dots,\ \mathbf{x}_K
  \label{eq:intro-traj}
\end{equation}
构成机器人的\emph{轨迹}。设地图由 $N$ 个\emph{路标}（landmark）组成，
\begin{equation}
  \mathbf{y}_1,\ \mathbf{y}_2,\ \dots,\ \mathbf{y}_N,
  \label{eq:intro-landmarks}
\end{equation}
每个时刻传感器测到其中一部分路标，得到观测数据。这样，“机器人携带传感器在环境中运动”就由两件事描述：\emph{什么是运动}（从 $k-1$ 到 $k$ 时刻 $\mathbf{x}$ 如何变化），\emph{什么是观测}（$k$ 时刻在 $\mathbf{x}_k$ 处看到路标 $\mathbf{y}_j$ 时得到什么数据）。

\subsection{运动方程}
\label{sec:intro-motion}
机器人携带测量自身运动的传感器（码盘、惯性传感器），其读数不一定直接是位置之差，可能是加速度、角速度等；或我们直接给它指令（“前进 1 米”“左转 $90^\circ$”“油门踩到底”“刹车”）。无论哪种，都可统一写成一个抽象的\emph{运动方程}（motion equation）\cite{gaoxiang2019slam14}：
\begin{equation}
  \mathbf{x}_k = f\!\left(\mathbf{x}_{k-1},\, \mathbf{u}_k,\, \mathbf{w}_k\right),
  \label{eq:intro-motion}
\end{equation}
其中 $\mathbf{u}_k$ 是运动传感器的读数或输入，$\mathbf{w}_k$ 是该过程中加入的\emph{噪声}，$f$ 是一个\emph{不指明具体作用方式}的一般函数——正因不指明，它才能指代任意运动传感器或输入，成为通用方程。

\begin{insight}
噪声 $\mathbf{w}_k$ 的存在，使 \cref{eq:intro-motion} 变成了一个\emph{随机模型}\cite{gaoxiang2019slam14}：即使你下达“前进 1 米”的命令，也不代表机器人真的前进了 1 米——它可能这次走了 $0.9$ 米、下次走了 $1.1$ 米，再一次因轮胎打滑根本没动。\emph{每次运动过程的噪声都是随机的}。若对噪声视而不见、只按指令推算位置，估计结果可能与实际相差十万八千里。这一句话，奠定了 SLAM “必须严格携带不确定性”的基调（贯穿 \cref{ch:slam_est}）。本书按统一约定取 $\mathbf{w}_k\sim\mathcal{N}(\mathbf{0},\boldsymbol{\Sigma}_w)$。
\end{insight}

\subsection{观测方程}
\label{sec:intro-obs}
与运动方程相对应，\emph{观测方程}（observation equation）描述机器人在 $\mathbf{x}_k$ 处看到路标 $\mathbf{y}_j$ 时产生观测数据 $\mathbf{z}_{k,j}$ 的过程\cite{gaoxiang2019slam14}：
\begin{equation}
  \mathbf{z}_{k,j} = h\!\left(\mathbf{y}_j,\, \mathbf{x}_k,\, \mathbf{v}_{k,j}\right),
  \label{eq:intro-obs}
\end{equation}
其中 $\mathbf{v}_{k,j}$ 是这次观测的噪声（本书取 $\mathbf{v}_{k,j}\sim\mathcal{N}(\mathbf{0},\boldsymbol{\Sigma}_v)$）。由于观测传感器形式更多，观测数据 $\mathbf{z}$ 与观测函数 $h$ 也有许多不同形式。

\subsection{两个具体参数化}
\label{sec:intro-param}
方程 \cref{eq:intro-motion,eq:intro-obs} 里的 $f,h$ 与 $\mathbf{x},\mathbf{y},\mathbf{z}$ 究竟长什么样，取决于真实运动与传感器的种类，存在若干\emph{参数化}（parameterization）方式\cite{gaoxiang2019slam14}。下面两个例子把抽象方程落到实处。

\begin{example}[平面运动 + 位移/角度增量输入]\label{ex:intro-planar}
设机器人在平面中运动，\emph{位姿}由两个位置和一个转角描述，$\mathbf{x}_k=[x_1,\,x_2,\,\theta]_k^\top$（$x_1,x_2$ 是两轴上的位置、$\theta$ 是转角）；输入指令是两个时间间隔内位置与转角的变化量，$\mathbf{u}_k=[\Delta x_1,\,\Delta x_2,\,\Delta\theta]_k^\top$。则运动方程 \cref{eq:intro-motion} 具体化为
\begin{equation}
  \begin{bmatrix} x_1\\ x_2\\ \theta \end{bmatrix}_{k}
  =\begin{bmatrix} x_1\\ x_2\\ \theta \end{bmatrix}_{k-1}
  +\begin{bmatrix} \Delta x_1\\ \Delta x_2\\ \Delta\theta \end{bmatrix}_{k}
  +\mathbf{w}_k.
  \label{eq:intro-planar-motion}
\end{equation}
这是一个\emph{简单的线性关系}。但并非所有输入都是位移/角度的变化量——“油门”“控制杆”的输入是速度或加速度，那时就会有更复杂的运动方程，需要做动力学分析。
\end{example}

\begin{note}
\textbf{严格性补注.}\;\cref{ex:intro-planar} 把转角 $\theta$ 直接相加，是最朴素的写法，没有处理角度的 $2\pi$ 环绕，本质上是小角度/局部近似；严格地说，平面位姿属于 $\SE(2)$，应在群上做\emph{乘法}而非分量相加。三维位姿则属 $\SE(3)$，其表示与运算是 \cref{ch:rigid_body,ch:lie} 的内容。
\end{note}

\begin{example}[二维激光观测一个 2D 路标]\label{ex:intro-laser}
设机器人携带二维激光传感器，观测一个 2D 路标点时能测到两个量：路标点与本体之间的\emph{距离} $r$ 和\emph{夹角} $\phi$。记路标 $\mathbf{y}_j=[y_1,\,y_2]_j^\top$、位姿 $\mathbf{x}_k=[x_1,\,x_2]_k^\top$、观测 $\mathbf{z}_{k,j}=[r_{k,j},\,\phi_{k,j}]^\top$，则观测方程 \cref{eq:intro-obs} 具体化为
\begin{equation}
  \begin{bmatrix} r_{k,j}\\[2pt] \phi_{k,j} \end{bmatrix}
  =\begin{bmatrix}
     \sqrt{(y_{1,j}-x_{1,k})^2+(y_{2,j}-x_{2,k})^2}\\[4pt]
     \arctan\dfrac{y_{2,j}-x_{2,k}}{y_{1,j}-x_{1,k}}
   \end{bmatrix}+\mathbf{v}_{k,j}.
  \label{eq:intro-laser-obs}
\end{equation}
而当传感器换成\emph{相机}（视觉 SLAM）时，观测方程就是“对路标点拍摄后得到图像中的像素”这一过程，它牵涉相机的成像模型，将在 \cref{ch:camera} 详述。
\end{example}

\begin{note}
\textbf{严格性补注.}\;\cref{eq:intro-laser-obs} 第二行的 $\arctan$ 严格应写成四象限的 \texttt{atan2}，以正确处理角度所在象限；原始写法是简化。噪声项 $\mathbf{v}_{k,j}\sim\mathcal{N}(\mathbf{0},\boldsymbol{\Sigma}_v)$、$\boldsymbol{\Sigma}_v\in\mathbb{R}^{2\times2}$。
\end{note}

\subsection{通用形式与“SLAM = 状态估计”}
\label{sec:intro-statevar}
保持通用性，把两方程取成抽象形式，SLAM 过程就可总结为两个基本方程\cite{gaoxiang2019slam14}：
\begin{equation}
  \left\{
  \begin{array}{ll}
    \mathbf{x}_k = f\!\left(\mathbf{x}_{k-1},\,\mathbf{u}_k,\,\mathbf{w}_k\right), & k=1,\dots,K,\\[4pt]
    \mathbf{z}_{k,j} = h\!\left(\mathbf{y}_j,\,\mathbf{x}_k,\,\mathbf{v}_{k,j}\right), & (k,j)\in\mathcal{O},
  \end{array}
  \right.
  \label{eq:intro-slam-general}
\end{equation}
其中 $\mathcal{O}$ 是一个\emph{集合}，记录在哪个时刻观察到了哪个路标——通常并非每个路标在每个时刻都能看到，单个时刻往往只看到一小部分。\cref{eq:intro-slam-general} 描述了最基本的 SLAM 问题：\textbf{当知道运动测量读数 $\mathbf{u}$ 与传感器读数 $\mathbf{z}$ 时，如何求解定位（估计 $\mathbf{x}$）与建图（估计 $\mathbf{y}$）？}

\begin{definition}[SLAM 作为状态估计问题]\label{def:intro-state-est}
将 \cref{eq:intro-slam-general} 视为一个\emph{状态估计问题}：\textbf{如何通过带有噪声的测量数据，估计内部的、隐藏着的状态变量}（机器人轨迹 $\{\mathbf{x}_k\}$ 与地图 $\{\mathbf{y}_j\}$）\cite{gaoxiang2019slam14}。这是后续整本书状态估计与优化部分（\cref{ch:slam_est,ch:nlopt,ch:eskf}）的统一出发点。
\end{definition}

\subsection{两种视角：递归贝叶斯（滤波）与 MAP（优化）}
\label{sec:intro-two-views}
\paragraph{为什么要两种视角。}
\cref{eq:intro-slam-general} 把 SLAM 写成了\emph{函数}形式；同一问题还能写成\emph{概率}形式，而正是这两种等价写法，分别孕育了求解 SLAM 的两大技术路线——滤波与优化。把它们并排放在一起，能让你在读后续章节时随时知道“现在站在哪条路上”。

\paragraph{视角一：递归贝叶斯（滤波的母方程）。}
Durrant-Whyte 与 Bailey\cite{durrantwhyte2006slam} 把运动与观测写成概率密度：状态转移是一阶\emph{马尔可夫过程}，下一状态只依赖前一状态与所施控制，
\begin{equation}
  P\!\left(\mathbf{x}_k\mid \mathbf{x}_{k-1},\mathbf{u}_k\right)\quad(\text{运动模型}),\qquad
  P\!\left(\mathbf{z}_k\mid \mathbf{x}_k,\mathbf{m}\right)\quad(\text{观测模型}),
  \label{eq:intro-dw-models}
\end{equation}
其中观测在“车辆位置与地图给定”后\emph{条件独立}。概率 SLAM 要求对所有 $k$ 计算联合后验 $P(\mathbf{x}_k,\mathbf{m}\mid Z_{0:k},U_{0:k},\mathbf{x}_0)$，它可由一对标准的\emph{预测—校正}递归实现\cite{durrantwhyte2006slam}：

\emph{时间更新（预测）}——用运动模型把上一时刻后验向前推一步：
\begin{equation}
  P\!\left(\mathbf{x}_k,\mathbf{m}\mid Z_{0:k-1},U_{0:k},\mathbf{x}_0\right)
  =\int P\!\left(\mathbf{x}_k\mid \mathbf{x}_{k-1},\mathbf{u}_k\right)
  P\!\left(\mathbf{x}_{k-1},\mathbf{m}\mid Z_{0:k-1},U_{0:k-1},\mathbf{x}_0\right)\,\mathrm{d}\mathbf{x}_{k-1}.
  \label{eq:intro-dw-predict}
\end{equation}

\emph{测量更新（校正）}——用新观测经贝叶斯法则修正预测：
\begin{equation}
  P\!\left(\mathbf{x}_k,\mathbf{m}\mid Z_{0:k},U_{0:k},\mathbf{x}_0\right)
  =\frac{P\!\left(\mathbf{z}_k\mid \mathbf{x}_k,\mathbf{m}\right)\,
         P\!\left(\mathbf{x}_k,\mathbf{m}\mid Z_{0:k-1},U_{0:k},\mathbf{x}_0\right)}
        {P\!\left(\mathbf{z}_k\mid Z_{0:k-1},U_{0:k}\right)}.
  \label{eq:intro-dw-update}
\end{equation}
当运动/观测模型取\emph{状态空间 + 加性高斯噪声}时，这对递归的最常见实现就是\emph{扩展卡尔曼滤波}（EKF）；若把运动模型描述为更一般非高斯分布的\emph{样本}，则导出 Rao-Blackwell 化的\emph{粒子滤波}（FastSLAM）\cite{durrantwhyte2006slam}。\cref{eq:intro-dw-predict,eq:intro-dw-update} 正是 \cref{ch:eskf} 滤波器一章的母方程；而 \cref{thm:intro-convergence} 的单调收敛，在 EKF-SLAM 中就表现为地图协方差行列式的单调下降。

\paragraph{视角二：MAP 与因子图（优化的母方程）。}
现代主流走的是另一条路——把 SLAM 写成一个\emph{最大后验估计}\cite{cadena2016past}。记待估变量 $X$（含整条轨迹与全部路标），测量集 $Z=\{\mathbf{z}_k\}$，每个测量写成 $\mathbf{z}_k=h_k(X_k)+\boldsymbol{\epsilon}_k$（$X_k\subseteq X$ 是它依赖的变量子集，$\boldsymbol{\epsilon}_k$ 是噪声）。由贝叶斯法则，
\begin{equation}
  X^\star=\arg\max_X p(X\mid Z)=\arg\max_X p(Z\mid X)\,p(X),
  \label{eq:intro-map}
\end{equation}
其中 $p(Z\mid X)$ 是似然、$p(X)$ 是先验（无先验知识时 $p(X)$ 为常数，MAP 退化为\emph{最大似然估计} MLE）。假设各测量独立，似然可\emph{因子分解}：
\begin{equation}
  X^\star=\arg\max_X\, p(X)\prod_{k=1}^{m} p(\mathbf{z}_k\mid X_k).
  \label{eq:intro-map-factor}
\end{equation}
\cref{eq:intro-map-factor} 可解读为在一张\emph{因子图}（factor graph）上的推断：变量是节点，每个 $p(\mathbf{z}_k\mid X_k)$ 与先验 $p(X)$ 是连接相关变量的\emph{因子}。若测量噪声为零均值高斯、信息矩阵为 $\boldsymbol{\Omega}_k$，则 $p(\mathbf{z}_k\mid X_k)\propto\exp(-\tfrac12\|h_k(X_k)-\mathbf{z}_k\|_{\boldsymbol{\Omega}_k}^2)$（记加权范数 $\|\mathbf{e}\|_{\boldsymbol{\Omega}}^2=\mathbf{e}^\top\boldsymbol{\Omega}\,\mathbf{e}$）；最大化后验等价于最小化负对数后验，于是 MAP 变成一个\emph{非线性最小二乘}问题：
\begin{equation}
  X^\star=\arg\min_X\ \sum_{k} \big\|h_k(X_k)-\mathbf{z}_k\big\|_{\boldsymbol{\Omega}_k}^2.
  \label{eq:intro-nls}
\end{equation}

\begin{insight}
\cref{eq:intro-nls} 是全书\emph{后端优化}的母方程，几乎所有现代 SLAM 系统都在解它的某个变体\cite{cadena2016past}。它与运动恢复结构（SfM）里的\emph{光束法平差}（Bundle Adjustment）同源，但 SLAM 有两点独特：因子\emph{不限于}投影几何，可纳入惯性、轮速、GPS、激光、光度等多种约束（这就是 \cref{sec:intro-backend} 说后端“不关心数据来自什么传感器”的数学体现）；且 SLAM 需\emph{增量}求解（机器人移动时新测量逐步到来）。求解时用高斯-牛顿或 Levenberg--Marquardt 逐次线性化，\emph{现代求解器的关键洞见是：正规方程的矩阵是稀疏的，稀疏结构由因子图的拓扑决定}——这正是算法分析时代（\cref{tab:intro-eras}）认识到的“稀疏性”，让 GTSAM、g2o、Ceres 等能在数秒内解出万级变量。其完整推导与实现见 \cref{ch:nlopt}。
\end{insight}

\begin{pitfall}[概念误区：滤波与优化是两个“不相干”的方法]
不要把递归贝叶斯（\cref{eq:intro-dw-predict,eq:intro-dw-update}）与 MAP（\cref{eq:intro-nls}）当成毫无关系的两套东西。二者求解的是\emph{同一个}联合后验，只是\emph{切法}不同：滤波\emph{边缘化}掉过去的位姿、用一个概率分布概括历史，每步只更新当前状态；优化（平滑）则\emph{保留}（部分）历史位姿，做全局调整。\textbf{在线性高斯情形，卡尔曼滤波与 MAP 给出完全相同的估计}\cite{cadena2016past}；分歧只在非线性时出现——滤波的线性化“一次性、不可回头”，优化可反复重新线性化，故通常更准。理解这一点，你才不会在 \cref{ch:eskf}（滤波）与 \cref{ch:nlopt}（优化）之间感到割裂。
\end{pitfall}

\subsection{按线性/高斯分类，引出滤波 vs 优化}
\label{sec:intro-est}
\paragraph{四象限分类。}
求解 \cref{eq:intro-slam-general} 用什么方法，取决于方程是否\emph{线性}、噪声是否服从\emph{高斯}分布\cite{gaoxiang2019slam14}，见 \cref{tab:intro-lgnlng}。

\begin{table}[H]\centering\small
\caption{按系统线性性与噪声分布分类的最优估计方法}
\label{tab:intro-lgnlng}
\begin{tabular}{p{0.26\textwidth} p{0.14\textwidth} p{0.12\textwidth} p{0.36\textwidth}}
\toprule
系统类型 & 方程 & 噪声 & 最优估计方法 \\
\midrule
线性高斯 (LG) & 线性 & 高斯 & \emph{卡尔曼滤波器}（KF）——最简单，给出\emph{无偏的最优估计} \\
非线性非高斯 (NLNG) & 非线性 & 非高斯 & \emph{扩展卡尔曼滤波}（EKF）与\emph{非线性优化}两大类（实际 SLAM 多属此类） \\
\bottomrule
\end{tabular}
\end{table}

\noindent EKF 的做法是\emph{在工作点处把系统线性化}，再以\emph{预测—更新}两步求解（即 \cref{eq:intro-dw-predict,eq:intro-dw-update} 的高斯实现），细节见 \cref{ch:eskf}。

\paragraph{历史与结论：优化为何胜出。}
回顾这条方法论的演变\cite{gaoxiang2019slam14}：直至 21 世纪早期，以 EKF 为主的滤波器方法在 SLAM 中\emph{占据主导}，最早的实时视觉 SLAM 系统就是基于 EKF 开发的；随后为克服 EKF 的两大缺点——\emph{线性化误差}与\emph{噪声高斯分布假设}——人们转向粒子滤波等其他滤波器，乃至非线性优化方法；时至今日，主流视觉 SLAM 使用以\emph{图优化}（Graph Optimization）为代表的优化技术做状态估计。本书的立场是：\textbf{优化技术已明显优于滤波器技术，只要计算资源允许，通常都偏向使用优化方法}\cite{gaoxiang2019slam14}。这一结论也得到独立研究的支持——Strasdat 等的“Visual SLAM: Why Filter?”指出，在路标数很大时，\emph{基于关键帧的光束法平差在每单位计算时间内给出的精度最高}，优于滤波；滤波只在算力极受限的小场景里仍有一席之地。两条路线分别在 \cref{ch:eskf}（滤波）与 \cref{ch:nlopt,ch:slam_est}（优化）展开。

\subsection{仍待澄清的问题（直接引出后续各章）}
\label{sec:intro-pending}
写到这里，几个问题被有意搁置，它们恰好牵出全书数学预备的主线\cite{gaoxiang2019slam14}：

\begin{enumerate}
  \item \textbf{机器人位置 $\mathbf{x}$ 到底是什么？}平面上可用“两坐标 + 一转角”，但机器人多数时候在三维空间运动：由 3 个轴上的平移 + 绕 3 个轴的旋转描述，共 \emph{6 个自由度}。能不能随便拿一个 $\mathbb{R}^6$ 向量来表示？\emph{事情没那么简单}——6 自由度位姿如何\emph{表达}、如何\emph{优化}，需要专门的篇幅，这是 \cref{ch:rigid_body}（旋转矩阵、变换矩阵、四元数、旋转向量）与 \cref{ch:lie}（李群与李代数、流形上的优化）的内容。
  \item \textbf{视觉 SLAM 的观测方程如何参数化？}即空间路标点如何投影成照片里的像素——这是相机成像模型，见 \cref{ch:camera}。
  \item \textbf{怎么求解上述方程？}从带噪测量估出隐藏状态，需要非线性优化的知识，见 \cref{ch:nlopt}。
\end{enumerate}

\noindent 可以看到，本章介绍的内容\emph{构成了本书的提要}——这些悬而未决的问题，逐一对应着后续的章节。下一节的全书路线图，会把这种“本章埋问题、后续章作答”的对应关系一次说全。

\begin{practice}
\begin{enumerate}
  \item 为 \cref{ex:intro-laser} 写出去掉噪声后的“无噪观测”，并验证：当机器人恰在路标正东方 $5\,$米处时，$r=5$、$\phi$ 取何值？说明为何此处必须用 \texttt{atan2} 而非 $\arctan$。
  \item 把 \cref{eq:intro-map}（MAP）在“无先验”时退化为 MLE，写出对应的 $\arg\max$ 表达式，并说明退化意味着什么假设。
  \item 用一句话分别概括滤波与优化“切”联合后验的不同方式，并说明二者在线性高斯下为何等价。
  \item 列出本章被搁置的三个问题，各指出它在哪一章解决（用 \cref{ch:rigid_body} 等标签思考）。
\end{enumerate}
\end{practice}

% ════════════════════════════════════════════════════════════════
\section{SLAM 的分类}
\label{sec:intro-class}

\paragraph{动机：给纷繁的 SLAM 系统一套坐标。}
读论文时你会被各种名字淹没：单目/双目/RGB-D、视觉/激光/视觉-惯性、滤波/优化、特征法/直接法……它们其实落在几个相互正交的分类维度上。把这套坐标建立起来，任何一个系统都能被你迅速定位。

\subsection{按传感器分类}
\label{sec:intro-by-sensor}
\begin{table}[H]\centering\small
\caption{按传感器对 SLAM 分类}
\label{tab:intro-class-sensor}
\begin{tabular}{p{0.16\textwidth} p{0.24\textwidth} p{0.48\textwidth}}
\toprule
维度 & 类别 & 说明 \\
\midrule
传感器主体 & 视觉 SLAM & 主要用相机；按相机再分单目/双目/RGB-D（本书主题） \\
 & 激光 SLAM & 主要用激光（2D/3D LiDAR）；相对成熟，有大量现成软件（\cref{ch:lidar_slam}） \\
 & 视觉-惯性 SLAM（VI-SLAM/VIN） & 相机 + IMU；IMU 提供高频运动与尺度；可视为禁用回环的简化 SLAM（\cref{ch:vio}） \\
相机类型 & 单目 Monocular & 结构简单成本低；有尺度不确定性、需平移才能算深度（\cref{thm:intro-scale}） \\
 & 双目 Stereo / 多目 & 已知基线估每像素深度；室内外皆可；计算量大 \\
 & 深度 RGB-D & 结构光/ToF 主动测距；省算力；主要室内 \\
\bottomrule
\end{tabular}
\end{table}

\noindent 现代系统还有更多组合：激光-惯性 SLAM（LIO）、激光-视觉-惯性 SLAM（LVI）、Radar SLAM、事件相机 SLAM 等。本书将分别在 \cref{ch:imu}（惯性测量）、\cref{ch:vio}（视觉-惯性）、\cref{ch:lidar_slam}（激光及激光-惯性）、\cref{ch:calib}（多传感器标定与可观性）展开这些融合方案。

\subsection{按方法分类}
\label{sec:intro-by-method}
\begin{table}[H]\centering\small
\caption{按方法对 SLAM 分类}
\label{tab:intro-class-method}
\begin{tabular}{p{0.18\textwidth} p{0.24\textwidth} p{0.46\textwidth}}
\toprule
维度 & 类别 & 说明 \\
\midrule
后端范式 & 滤波（Filtering） & EKF/UKF/粒子滤波；边缘化过去位姿、用分布概括历史；21 世纪早期主导（\cref{ch:eskf}） \\
 & 优化/平滑（Optimization/Smoothing） & 图优化/因子图/BA/位姿图；保留（部分）历史位姿做全局优化；当代主流（\cref{ch:nlopt}） \\
系统线性/噪声 & 线性高斯 (LG) & KF 给无偏最优估计 \\
 & 非线性非高斯 (NLNG) & EKF / 非线性优化 \\
视觉前端范式 & 特征点法（feature-based/indirect） & 提取-匹配特征点、最小化\emph{重投影误差}；显式数据关联、收敛域大、较稳健、成熟（\cref{ch:vo}） \\
 & 直接法（direct） & 直接最小化\emph{光度误差}、不需特征匹配；与光流相关；常 coarse-to-fine（\cref{ch:vo}） \\
 & 半直接/混合 & 介于二者之间（如 SVO） \\
\bottomrule
\end{tabular}
\end{table}

\begin{insight}
\textbf{“为何不用滤波？”的权威结论}：Strasdat 等指出，滤波\emph{边缘化}过去位姿、用概率分布概括历史，而基于关键帧的优化\emph{保留}全局 BA 的结构、只选少量过去帧；在大多数现代应用中，\emph{关键帧优化在每单位计算时间内给出的精度最高}，路标数很大时优于滤波。滤波只在算力极低的系统里仍有用武之地。这与本书 \cref{sec:intro-est} 的结论（优化已明显优于滤波）一致，也解释了为何本书把优化作为后端主线、滤波作为重要补充。
\end{insight}

% ════════════════════════════════════════════════════════════════
\section{现代代表系统概览}
\label{sec:intro-systems}

\paragraph{动机：给方法论装上真实的坐标。}
前两节给了分类的“坐标轴”，这一节给出落在坐标上的“点”——若干里程碑系统。它们既是历史，也是你日后实现与对比时的参照。先看两个奠基者，再看三类现代主流。

\paragraph{两个奠基系统。}
\emph{MonoSLAM}（Davison, 2007）是\textbf{首个实时单目视觉 SLAM} 系统，用 EKF 维护相机与少量特征点的联合状态；其复杂度随路标数呈 $O(N^3)$，只适合几百个点的小场景——这正是 EKF-SLAM 计算量随路标数平方/立方增长的实证。\emph{PTAM}（Klein \& Murray, 2007）做了两件影响深远的事：\emph{首次用非线性优化（BA）替代滤波}做后端，并\emph{首次把跟踪（tracking）与建图（mapping）并行化}为双线程。自 PTAM 起，把 SfM 的光束法平差引入视觉 SLAM 成为主流，也奠定了 \cref{sec:intro-frontback} 所说的并行流水线范式。

\paragraph{三类现代主流。}
\cref{tab:intro-modern} 给出对应纯视觉、视觉-惯性、激光-惯性三条线的代表系统。

\begin{table}[H]\centering\small
\caption{三类现代主流 SLAM 系统（按本书后续章节归属）}
\label{tab:intro-modern}
\begin{tabular}{p{0.16\textwidth} p{0.18\textwidth} p{0.50\textwidth}}
\toprule
系统 & 传感器 & 核心特征 \\
\midrule
ORB-SLAM 家族\cite{murartal2015orbslam} & 单目 / 双目 / RGB-D (+IMU) & 特征点法标杆；\emph{三线程并行}（tracking / local mapping / loop closing）；\emph{全任务共用同一 ORB 特征}；用 DBoW2 二进制词袋做快速地点识别；回环用 Sim(3) 修正单目尺度漂移；ORB-SLAM3 进一步做\emph{完全基于 MAP 的紧耦合视觉-惯性 + 多地图 ATLAS}（\cref{ch:vo,ch:loop,ch:vio}） \\
VINS-Mono\cite{qin2018vins} & 单目 + 低成本 IMU & 优化型 VIO 标杆；度量 6-DOF 状态估计的\emph{最小传感器组合}；\emph{滑动窗口紧耦合}非线性优化融合 IMU 预积分与特征观测；含鲁棒初始化与失败恢复、回环 + 重定位、\emph{4-DOF 位姿图优化}（重力使 roll/pitch 可观）（\cref{ch:imu,ch:vio}） \\
LIO-SAM\cite{shan2020liosam} & 激光 + IMU & 因子图平滑型 LIO 标杆；把激光-惯性里程计建在\emph{因子图}上，纳入四类因子（IMU 预积分、激光里程计、GPS、回环）；用 IMU 预积分做点云去畸变并提供初值，激光解又用于估 IMU 偏置；\emph{边缘化}旧扫描以保实时（\cref{ch:lidar_slam}） \\
\bottomrule
\end{tabular}
\end{table}

\noindent 这三者恰好对应三种现代主流：ORB-SLAM = 纯视觉（特征点法 + 优化 + 多地图），VINS-Mono = 视觉-惯性（紧耦合滑窗优化），LIO-SAM = 激光-惯性（因子图平滑）。它们都用到本章 \cref{eq:intro-nls} 的优化母方程与因子图思想，区别只在喂进去的是哪种因子——再一次印证了后端“只面对数据”的统一性。

% ════════════════════════════════════════════════════════════════
\section{前沿与展望：SLAM 还能走多远}
\label{sec:intro-frontier}

\paragraph{动机：经典框架的边界之外。}
本章反复强调过经典视觉 SLAM 框架（\cref{sec:intro-pipeline}）有一个隐含前提——\emph{静态、刚体、光照稳定、无人干扰}的环境（\cref{tab:intro-modules} 之后已点明）。现实世界恰恰常常\emph{不}满足这些前提：场景里有走动的人、行驶的车、随风摆动的树，甚至会形变的柔软物体；任务规模又可能大到一台机器人吃不下，需要多台协同。本节给出几条把 SLAM 推出经典边界的活跃方向。它们多处于鲁棒感知时代（\cref{tab:intro-eras}）“追求鲁棒性、追求高层理解”的诉求之下，是\emph{展望性}的概览——本书不会逐一展开其全部细节，而是指明问题、给出代表工作与它们各自归属的后续专章。除下面两条外，本章前文已把另两条重要方向落到正文：\emph{视觉-惯性融合}（用 IMU 抵御快速运动与动态干扰）见 \cref{ch:vio}，\emph{语义与深度学习}（让地图带上物体标签、用网络做地点识别）见 \cref{ch:dense,ch:loop}。

\subsection{动态与可变形场景 SLAM}
\label{sec:intro-dynamic}
经典框架把环境当作\emph{刚性且静止}的“背景板”，凡是运动的像素都被当成相机自身运动的线索。可一旦场景里有\emph{独立运动}的物体，这个假设就破裂了：动态物体上的特征会把错误的运动“注入”位姿估计，轻则降低精度、重则跟踪发散。如何让 SLAM 在动态世界里仍然稳健，是当前的活跃方向，思路大体分两路\cite{cadena2016past}。

\emph{一是把动态物当作干扰，检测并剔除}。代表工作 DynaSLAM\cite{bescos2018dynaslam} 在 ORB-SLAM2 之上加一个动态物体模块：用 Mask R-CNN 做语义分割、再结合多视图几何识别出运动物体，把落在其上的特征\emph{剔除}，只用静态背景做定位，并\emph{修补}（inpaint）被动态物遮挡的背景区域，得到一张“干净”的静态地图。这一路与语义 SLAM（\cref{ch:dense}）天然耦合——“哪些像素属于可能移动的物体”往往要靠语义来回答。

\emph{二是正面建模运动甚至形变，分别跟踪}。代表工作 Co-Fusion\cite{runz2017cofusion} 实时地把 RGB-D 场景\emph{分割}成若干物体，为背景与每个独立运动的物体各自维护一个三维模型、分别跟踪与融合。更进一步，若物体本身会\emph{非刚性形变}（人脸、衣物、器官），刚体假设彻底失效——DynamicFusion\cite{newcombe2015dynamicfusion} 是这方面的里程碑：它无须任何先验模板，实时估计一个稠密的\emph{体素形变场}（warp field），把每一帧非刚性地“拉回”到一个标准模型上，从而重建并跟踪可变形场景。这类稠密重建与本书的稠密建图（\cref{ch:dense}）一脉相承，只是把“静态体素”推广成了“随时间形变的体素”。

\begin{insight}
动态 SLAM 把本章“鸡生蛋”循环（\cref{sec:intro-chicken-egg}）又拧紧了一圈：经典 SLAM 要同时估\emph{位姿}与\emph{静态地图}，动态场景下还要同时判断\emph{每个观测到底属于静态背景还是某个运动物体}——这本质上是一个更难的\emph{数据关联}问题（\cref{sec:intro-frontback}）。语义信息之所以在这里举足轻重，正因为它为“这块像素会不会动”提供了先验。
\end{insight}

\subsection{多机器人与协作 SLAM}
\label{sec:intro-multirobot}
本章至此默认只有\emph{一台}机器人。但许多任务——大范围测绘、灾后搜救、仓储调度——用\emph{一队}机器人协同完成会快得多、也更鲁棒。\emph{多机器人}（协作）SLAM 研究的就是：多台机器人如何各自建图、又把彼此的地图与轨迹\emph{拼成一张全局一致的大图}\cite{cadena2016past}。它在单机 SLAM 之上引入两个新问题：一是\emph{机器人间回环}（inter-robot loop closure）——识别出“甲走过的地方乙也走过”，这是把多张子地图对齐的关键约束，本质上仍是\cref{sec:intro-loop} 回环检测的跨机器人版本；二是\emph{通信与算力的分配}——地图与约束在哪里融合。

按融合的地点，方案分两大类。\emph{集中式}（centralized）：每台机器人只跑轻量前端（VO），把数据汇给一台\emph{中央服务器}做地图合并、地点识别与全局优化——代表工作 CCM-SLAM\cite{schmuck2019ccmslam} 即为机器人小队设计了这样一套能容忍通信延迟与丢包的集中式协作单目 SLAM。\emph{分布式}（distributed）：没有中心节点，机器人两两通信、各自维护并优化地图。分布式的最大难点是\emph{错误的机器人间回环}（外观相似实为不同地点，即感知混叠，\cref{sec:intro-need}）会污染所有人的地图——DOOR-SLAM\cite{lajoie2020doorslam} 用\emph{两两一致测量集}（pairwise consistent measurement set, PCM）筛除离群回环，实现了在线、抗外点的分布式协作 SLAM。再进一步，Kimera-Multi\cite{tian2022kimeramulti} 把分布式协作与度量-语义建图结合：仅靠点对点通信即可做分布式地点识别与鲁棒位姿图优化，最终拼出一张全局一致、带语义标注的三维网格地图。

\begin{insight}
多机器人 SLAM 仍然在解本章 \cref{eq:intro-nls} 的那条优化母方程——只是因子图（\cref{sec:intro-two-views}）的变量来自\emph{多台}机器人、其中一类特殊因子是“机器人间回环”。集中式与分布式之争，本质是\emph{这张大因子图在哪里、由谁来优化}：集中到服务器一次解完，还是切成子图、由各机器人协同迭代。无论哪种，抗感知混叠的\emph{鲁棒回环}都是命门——这与单机 SLAM 把回环当成可靠性基石（\cref{sec:intro-need}）的判断完全一致。
\end{insight}

% ════════════════════════════════════════════════════════════════
\section{实践：第一个 SLAM 工程}
\label{sec:intro-code}

\paragraph{动机：把环境立起来，后面才能动手。}
理论之外，本书每一章都有可运行的代码。它们以 \emph{Linux 上的 C++} 为主——大量 SLAM 库只对 Linux 提供较好支持（推荐 Ubuntu）。本节用最小的例子，把“源文件 $\to$ 编译 $\to$ 库 $\to$ 链接 $\to$ 可执行程序”这条工程链条讲透；之后各章就默认你已掌握这套 CMake 流程，不再重复。

\subsection{Hello SLAM：从 g++ 到第一个可执行程序}
\label{sec:intro-hello}
在 Linux 中，\emph{程序就是一个具有执行权限的文件}（脚本或二进制，不限后缀名）——常用的 \texttt{cd}、\texttt{ls} 就是 \texttt{/bin} 下的可执行文件。我们写下第一个程序：

\begin{codebox}[C++]{helloSLAM.cpp：第一个程序}
#include <iostream>
using namespace std;

int main(int argc, char **argv) {
    cout << "Hello SLAM!" << endl;
    return 0;
}
\end{codebox}

\noindent 用 C++ 编译器 g++ 把它编译成可执行程序：

\begin{codebox}[bash]{用 g++ 直接编译并运行}
# compile; if it succeeds there is no output at all.
# if you see "command not found", install the compiler first:
#   sudo apt-get install g++
g++ helloSLAM.cpp
# g++ names the product a.out by default (with executable permission); run it:
./a.out
# terminal prints: Hello SLAM!
\end{codebox}

\noindent 这一步用了大量默认参数、几乎省略所有中间步骤。当工程变大时，g++ 命令会变得无法维护——于是引入 CMake。

\subsection{用 CMake 管理工程}
\label{sec:intro-cmake}
\paragraph{为什么需要 CMake。}
理论上任意 C++ 程序都能用 g++ 编译；但当工程规模变大（许多文件夹与源文件、十几个类、复杂依赖、一部分编成可执行程序、一部分编成库），仅靠 g++ 需要输入大量编译指令，过程异常烦琐。历史上工程师用 \texttt{makefile} 自动编译，而 \texttt{cmake} 比它更方便、在工程上被广泛使用（本书后面用到的大多数库都用 CMake 管理）。CMake 的工作机制是：用 \texttt{cmake} 命令生成一个 \texttt{Makefile}，再用 \texttt{make} 据其编译整个工程。最小工程只需一个 \texttt{CMakeLists.txt}：

\begin{codebox}[text]{CMakeLists.txt：最小工程}
# declare the minimum required cmake version
cmake_minimum_required( VERSION 2.8 )
# declare a cmake project
project( HelloSLAM )
# add an executable.  syntax: add_executable( name source-files )
add_executable( helloSLAM helloSLAM.cpp )
\end{codebox}

\noindent 推荐用\emph{外部构建}（out-of-source build）：把中间文件隔离到一个 \texttt{build} 目录，发布时删掉它即可，不污染源码。

\begin{codebox}[bash]{out-of-source 构建（推荐做法）}
mkdir build
cd build
cmake ..     # configure against the parent dir (where the sources live);
             # all intermediate files are generated inside build/
make         # make actually calls g++ and produces the executable helloSLAM
./helloSLAM  # terminal prints: Hello SLAM!
\end{codebox}

\noindent CMake 多了 \texttt{cmake}、\texttt{make} 两步，却把工程管理从“输入一长串 g++ 命令”变成“维护若干直观的 \texttt{CMakeLists.txt}”，显著降低维护难度——新增一个可执行文件只需加一行 \texttt{add\_executable}，CMake 自动解决依赖关系。

\subsection{库、头文件与链接}
\label{sec:intro-lib}
\paragraph{库是什么。}
并非所有代码都编成可执行文件——\emph{只有带 \texttt{main} 函数的文件才生成可执行程序}；另一些代码只想打包供其他程序调用，这就是\emph{库}（Library）。库往往是许多算法的集合，比如 OpenCV 提供计算机视觉算法、Eigen 提供矩阵代数。先写一个没有 \texttt{main} 的库源文件：

\begin{codebox}[C++]{libHelloSLAM.cpp：一个库源文件（无 main）}
// this is a library source file
#include <iostream>
using namespace std;

void printHello() {
    cout << "Hello SLAM" << endl;   // a function provided by the library
}
\end{codebox}

\noindent 在 \texttt{CMakeLists.txt} 里可把它编成\emph{静态库}或\emph{共享库}（见 \cref{tab:intro-libtype}）：

\begin{codebox}[text]{把源文件编成库}
# static library -> produces libhello.a
add_library( hello libHelloSLAM.cpp )
# shared library -> produces libhello_shared.so
add_library( hello_shared SHARED libHelloSLAM.cpp )
\end{codebox}

\begin{table}[H]\centering\small
\caption{静态库与共享库}
\label{tab:intro-libtype}
\begin{tabular}{p{0.16\textwidth} p{0.10\textwidth} p{0.36\textwidth} p{0.18\textwidth}}
\toprule
库类型 & 后缀 & 调用行为 & 空间 \\
\midrule
静态库 static & \texttt{.a} & 每次被调用都生成一个副本 & 占空间多 \\
共享库 shared & \texttt{.so} & 全局只有一个副本 & 更省空间 \\
\bottomrule
\end{tabular}
\end{table}

\paragraph{头文件：库的“使用说明”。}
光有 \texttt{.a}/\texttt{.so} 库文件，调用者并不知道里面有哪些函数、怎么调用——因为库是编译好的二进制。要给别人（或自己）用，须提供\emph{头文件}说明库里有什么。使用者只要拿到\emph{头文件 + 库文件}即可调用：

\begin{codebox}[C++]{libHelloSLAM.h：库的头文件}
#ifndef LIBHELLOSLAM_H_
#define LIBHELLOSLAM_H_
// the include guard above prevents this header from being included twice
void printHello();   // declare the function provided by the library
#endif
\end{codebox}

\begin{codebox}[C++]{useHello.cpp：调用库的可执行程序}
#include "libHelloSLAM.h"   // refer to the library header to know how to call
int main(int argc, char **argv) {
    printHello();
    return 0;
}
\end{codebox}

\noindent 最后，在 \texttt{CMakeLists.txt} 里生成可执行程序并\emph{链接}到库：

\begin{codebox}[text]{生成可执行程序并链接到库}
add_executable( useHello useHello.cpp )
target_link_libraries( useHello hello_shared )   # link useHello against the lib
\end{codebox}

\begin{insight}
\textbf{库使用三要点}（贯穿全书每一个工程）：(1) 程序代码由\emph{头文件}和\emph{源文件}组成；(2) 带 \texttt{main} 的源文件编成\emph{可执行程序}，其余编成\emph{库文件}；(3) 可执行程序若要调用库里的函数，须\emph{参考库的头文件}（以明白调用格式），并把可执行程序\emph{链接}到库文件上\cite{gaoxiang2019slam14}。对他人提供的库（OpenCV、Eigen、Sophus、g2o、Ceres、GTSAM……），用法完全相同。
\end{insight}

\begin{pitfall}[陷阱：忘了链接库会怎样]
如果你 \texttt{\#include} 了库的头文件、调用了库函数，却忘了在 \texttt{CMakeLists.txt} 里写 \texttt{target\_link\_libraries}，编译\emph{不会}在编译阶段报错——错误发生在\emph{链接}阶段，形如 \texttt{undefined reference to printHello()}（未定义引用）。原因是：头文件只给了函数\emph{声明}，链接器找不到函数的\emph{实现}（它在库里）。记住这个错误样子，能帮你日后省下大量排错时间。
\end{pitfall}

\paragraph{把库装到本地，再用 \texttt{find\_package} 复用。}
上面的库与调用者在\emph{同一个}工程里，所以直接写库名即可链接。但实际工程中，库往往是\emph{别处}编译、\emph{安装}到系统后被许多工程复用的（OpenCV、Eigen、Sophus、g2o、Ceres、GTSAM 莫不如此）。其约定是：库作者用 \texttt{install} 规则把头文件与 \texttt{.a}/\texttt{.so} 装到本地标准目录（如 \texttt{/usr/local}），并附带一个 \emph{CMake 配置文件}（\texttt{xxxConfig.cmake}）说明库装在哪、叫什么；使用方则在 \texttt{CMakeLists.txt} 里用 \texttt{find\_package} 自动找到它，再像上面一样链接——见 \cref{tab:intro-findpkg}。这套机制把“头文件 + 库文件”这对凭据从\emph{手工指定路径}升级为\emph{按名查找}，是后续每一章引第三方库的通用写法。

\begin{table}[H]\centering\small
\caption{自带库链接 vs.\ \texttt{find\_package} 复用已装库}
\label{tab:intro-findpkg}
\begin{tabular}{p{0.18\textwidth} p{0.36\textwidth} p{0.38\textwidth}}
\toprule
环节 & 同工程内自带库 & 复用已安装的库（\texttt{find\_package}）\\
\midrule
库的来源 & 本工程内 \texttt{add\_library} 现编 & 别处编译、经 \texttt{install} 装到本地 \\
找到库 & 直接写库目标名 & \texttt{find\_package( Xxx REQUIRED )} 按名查找 \\
包含头文件 & 同目录可见 & \texttt{target\_include\_directories(\dots \$\{Xxx\_INCLUDE\_DIRS\})} \\
链接 & \texttt{target\_link\_libraries( app hello\_shared )} & \texttt{target\_link\_libraries( app \$\{Xxx\_LIBRARIES\})} \\
\bottomrule
\end{tabular}
\end{table}

\paragraph{IDE 与调试。}
工程文件一多，纯文本编辑器跳转、查声明就很烦，故推荐用 IDE（集成开发环境，如 KDevelop、CLion、VS Code）：它们支持 CMake 工程，提供高亮、跳转、补全与\emph{断点调试}。调试的关键是把工程设为 \emph{Debug} 模式（在 \texttt{CMakeLists.txt} 中 \texttt{set( CMAKE\_BUILD\_TYPE "Debug" )}，运行较慢但带调试信息），随后即可设断点、单步执行、查看中间变量的值——程序崩溃时用它定位出错位置极其有效。具体的最小调试流程见下面的 \cref{prac:intro-debug}（与 IDE 品牌无关，跑通一次即可）。从下一章起，本书默认你已掌握“新建 build、调用 cmake/make、用 IDE 调试”这套流程，不再重复。

\begin{practice}[最小调试流程：四步定位一个变量]\label{prac:intro-debug}
本书不复刻某款 IDE 的具体界面（KDevelop、CLion、VS Code 的菜单各异），只给出一套\emph{通用}的最小调试流程，对任意支持 CMake 的 IDE 都适用\cite{gaoxiang2019slam14}。以 \cref{sec:intro-lib} 的 \texttt{useHello} 为例：
\begin{enumerate}
  \item \textbf{设 Debug 模式}：在 \texttt{CMakeLists.txt} 中写 \texttt{set( CMAKE\_BUILD\_TYPE "Debug" )}（等价于给 g++ 加 \texttt{-g}、关掉 \texttt{-O2} 优化）。Debug 比 Release 慢，但带调试信息，断点和变量查看才有效。
  \item \textbf{设断点}：用 IDE 打开工程，在 \texttt{useHello.cpp} 里 \texttt{printHello();} 那一行左侧单击，放下一个断点。
  \item \textbf{以 Debug 方式启动}：指定要调试的可执行目标（这里是 \texttt{useHello}），点 “Debug/调试”（而非普通 “Run/运行”）启动；程序会\emph{停在}断点处等待。
  \item \textbf{单步执行、查看变量}：用单步（多为 F10，逐行不进函数）、单步进入（多为 F11，跟进 \texttt{printHello} 内部）控制执行，并在变量面板里观察当前各变量的值；看完点继续/停止结束调试。
\end{enumerate}
日后若程序在\emph{运行阶段}崩溃，就用同一套流程：在可疑行设断点、单步逼近，看哪一步的变量变成了非预期值——这比反复加打印语句高效得多。
\end{practice}

\begin{practice}
\begin{enumerate}
  \item 把 \cref{sec:intro-hello} 的程序敲一遍并用 g++ 编译运行；再查 g++ 手册，用哪个参数能把输出文件名从 \texttt{a.out} 改成 \texttt{hello}？（提示：\texttt{-o}。另查 \texttt{-std=c++11}、\texttt{-g}、\texttt{-O2}、\texttt{-Wall} 的含义。）
  \item 用 \texttt{build} 目录的 out-of-source 方式编译 \cref{sec:intro-cmake} 的工程，确认源码目录干净。
  \item 把 \cref{sec:intro-lib} 的 \texttt{target\_link\_libraries} 一行删掉再编译，记录报错信息，验证 \cref{sec:intro-lib} 陷阱说的“链接阶段 undefined reference”。
  \item 在代码里故意制造一个语法错误，看看 g++ 报什么；试着只凭报错信息定位它。
  \item \textbf{跑通一次断点调试（任选一款 IDE）}：照 \cref{prac:intro-debug} 的四步，在 \texttt{useHello.cpp} 调用 \texttt{printHello()} 那一行设断点、以 Debug 方式启动、单步进入函数、观察输出过程。把这套“设 Debug 模式 $\to$ 断点 $\to$ 单步 $\to$ 查变量”流程跑顺，之后各章遇到崩溃就能自行定位。
  \item \textbf{安装并复用自写库}：给 \cref{sec:intro-lib} 的工程加上 \texttt{install} 规则，把 \texttt{libhello\_shared.so} 与 \texttt{libHelloSLAM.h} 装到本地（如 \texttt{make install} 到 \texttt{/usr/local}）；再\emph{另建}一个全新工程，仅用 \cref{tab:intro-findpkg} 右列的 \texttt{find\_package}/\texttt{target\_link\_libraries} 找到并调用 \texttt{printHello()}。体会“一次安装、处处复用”——这正是后续引 OpenCV、Eigen、g2o 等第三方库的方式\cite{gaoxiang2019slam14}。
\end{enumerate}
\end{practice}

% ════════════════════════════════════════════════════════════════
\section{全书路线图：本章如何引出后续每一章}
\label{sec:intro-roadmap}

\paragraph{本章 = 全书地图。}
本章是全书的提要：它概括地介绍了一个视觉 SLAM 系统的结构，作为后续内容的大纲\cite{gaoxiang2019slam14}。下面把本章埋下的每一个问题、框架里的每一个模块，都对应到它将被解决的章节，见 \cref{tab:intro-roadmap}。这张表就是你的导航：读到任何一章卡住时，回到这里看它的上下文。

\begin{table}[H]\centering\small
\caption{本章埋下的问题 / 模块 $\to$ 后续章节}
\label{tab:intro-roadmap}
\begin{tabular}{p{0.40\textwidth} p{0.20\textwidth} p{0.30\textwidth}}
\toprule
本章埋下的问题 / 模块 & 去向章节 & 在那一章解决什么 \\
\midrule
6 自由度位姿如何\emph{表达}（\cref{sec:intro-pending}） & \cref{ch:rigid_body} & 旋转矩阵、变换矩阵、四元数、旋转向量、欧拉角 \\
6 自由度位姿如何\emph{优化}（\cref{sec:intro-pending}） & \cref{ch:lie} & 李群与李代数、指数/对数、扰动求导、流形优化 \\
观测方程如何参数化：点 $\to$ 像素（\cref{ex:intro-laser}） & \cref{ch:camera} & 针孔模型、内参、畸变、双目/RGB-D 模型 \\
怎么求解非线性最小二乘（\cref{eq:intro-nls}） & \cref{ch:nlopt} & 最小二乘、高斯-牛顿、LM、BA 与图优化、稀疏性 \\
状态估计的统一理论（\cref{def:intro-state-est}） & \cref{ch:slam_est} & MAP/MLE、不确定性、因子图、位姿图 \\
前端 VO：相邻帧运动估计（\cref{sec:intro-vo}） & \cref{ch:vo} & 特征点法（对极几何、PnP、ICP、三角化）与直接法/光流 \\
回环检测：图像相似性（\cref{sec:intro-loop}） & \cref{ch:loop} & 词袋模型、相似性评分、感知混叠、回环验证 \\
建图：稠密重建（\cref{sec:intro-mapping}） & \cref{ch:dense} & 单目/RGB-D 稠密建图、八叉树、TSDF \\
地图的点云表示与处理（\cref{sec:intro-mapping}） & \cref{ch:pointcloud} & 点云滤波、配准、表面重建 \\
滤波：预测-校正（\cref{eq:intro-dw-predict,eq:intro-dw-update}） & \cref{ch:eskf} & KF/EKF/ESKF、滑动窗口滤波 \\
惯性测量与预积分（\cref{tab:intro-class-sensor}） & \cref{ch:imu} & IMU 模型、噪声、预积分 \\
视觉-惯性融合（\cref{sec:intro-systems}） & \cref{ch:vio} & 紧耦合 VIO、初始化、可观性 \\
激光及激光-惯性 SLAM（\cref{tab:intro-class-sensor}） & \cref{ch:lidar_slam} & 激光里程计、LIO、因子图平滑 \\
多传感器标定与可观性（\cref{sec:intro-by-sensor}） & \cref{ch:calib} & 外参/时间标定、可观性与退化 \\
基于地图的导航与路径规划（\cref{sec:intro-mapping}） & \cref{ch:planning} & A*、D* 等路径搜索与运动规划 \\
完整 SLAM 系统的搭建（\cref{fig:intro-pipeline}） & \cref{ch:slam_system} & 把前端+后端+回环+建图集成为可运行系统 \\
记号与全书约定（本章记号注） & \nameref{ch:notation} & 全书统一记号、坐标系、扰动约定 \\
\bottomrule
\end{tabular}
\end{table}

\paragraph{两条隐藏主线。}
通读全书时，留意本章数学表述（\cref{sec:intro-math}）分出的两条主线：\emph{滤波主线}以递归贝叶斯的预测-校正（\cref{eq:intro-dw-predict,eq:intro-dw-update}）为母方程，落在 \cref{ch:eskf}；\emph{优化主线}以 MAP/因子图的非线性最小二乘（\cref{eq:intro-nls}）为母方程，落在 \cref{ch:nlopt,ch:slam_est}，并把回环（\cref{ch:loop}）当作“长期数据关联”因子自然纳入。本章的分类（\cref{sec:intro-class}）与现代系统（\cref{sec:intro-systems}）则是这两条主线在“模块/融合”层的预告。带着这张地图，你就不会在后续任何一章迷路。

% ════════════════════════════════════════════════════════════════
\section*{本章常见误解汇总}
\begin{table}[H]\centering\small
\begin{tabular}{p{0.46\textwidth} p{0.46\textwidth}}
\toprule
\textbf{常见误解} & \textbf{正确理解} \\
\midrule
先建图再定位、或先定位再建图即可 & 二者互为前提（“鸡生蛋”），须\emph{联合}估计（\cref{sec:intro-chicken-egg}）\\
联合后验可拆成定位后验 $\times$ 建图后验 & 观测同依赖位姿与地图，拆分得不一致估计（\cref{eq:intro-bad-partition}）\\
把相邻帧运动一段段串起来就有可信轨迹 & 误差累积成漂移，须靠回环 + 后端校正（\cref{ex:intro-drift}）\\
单目相机能直接测出真实尺度 & 存在尺度不确定性，轨迹/地图差一个未知因子（\cref{thm:intro-scale}）\\
后端要知道数据来自哪种传感器 & 后端只面对数据/因子，传感器差异在前端吸收（\cref{sec:intro-backend}）\\
滤波与优化是不相干的两套方法 & 解同一后验，线性高斯下等价；切法不同（\cref{sec:intro-two-views}）\\
6 自由度位姿可随便用一个 $\mathbb{R}^6$ 向量表示 & 位姿在流形上，表达/优化需 \cref{ch:rigid_body,ch:lie}\\
忘了链接库会在编译阶段报错 & 报错在\emph{链接}阶段：undefined reference（\cref{sec:intro-lib} 陷阱）\\
\bottomrule
\end{tabular}
\end{table}

% ════════════════════════════════════════════════════════════════
\section*{本章小结}

\begin{table}[H]\centering\small
\caption{符号表}
\begin{tabular}{lll}
\toprule
符号 & 含义 & 首次出现 \\
\midrule
$\mathbf{x}_k$ & $k$ 时刻机器人位姿/状态 & \cref{eq:intro-traj} \\
$\mathbf{y}_j$ & 第 $j$ 个路标/地图点；$\mathbf{m}$ 为地图整体 & \cref{eq:intro-landmarks} \\
$\mathbf{u}_k$ & 运动输入/控制 & \cref{eq:intro-motion} \\
$\mathbf{z}_{k,j}$ & $k$ 时刻对 $\mathbf{y}_j$ 的观测 & \cref{eq:intro-obs} \\
$\mathbf{w}_k,\mathbf{v}_{k,j}$ & 运动/观测噪声，$\sim\mathcal{N}(\mathbf{0},\boldsymbol{\Sigma}_{w/v})$ & \cref{eq:intro-motion,eq:intro-obs} \\
$f,h$ & 运动函数 / 观测函数 & \cref{eq:intro-motion,eq:intro-obs} \\
$\mathcal{O}$ & 观测索引集合（何时见何路标） & \cref{eq:intro-slam-general} \\
$\boldsymbol{\Sigma},\boldsymbol{\Omega}$ & 协方差 / 信息矩阵（$\boldsymbol{\Omega}=\boldsymbol{\Sigma}^{-1}$） & \cref{eq:intro-nls} \\
$X,Z$ & 全体待估变量 / 全体测量 & \cref{eq:intro-map} \\
$\|\mathbf{e}\|_{\boldsymbol{\Omega}}^2$ & 加权（马氏）范数 $\mathbf{e}^\top\boldsymbol{\Omega}\mathbf{e}$ & \cref{eq:intro-nls} \\
\bottomrule
\end{tabular}
\end{table}

\begin{table}[H]\centering\small
\caption{速查表}
\begin{tabular}{lll}
\toprule
要点 & 一句话说明 & 对应 \\
\midrule
运动/观测方程 & SLAM 两条母方程 & \cref{eq:intro-slam-general} \\
SLAM = 状态估计 & 由带噪测量估隐藏状态 & \cref{def:intro-state-est} \\
滤波母方程 & 递归贝叶斯预测-校正 & \cref{eq:intro-dw-predict,eq:intro-dw-update} \\
优化母方程 & MAP $\to$ 非线性最小二乘 & \cref{eq:intro-nls} \\
五模块框架 & 传感器/VO/后端/回环/建图 & \cref{fig:intro-pipeline} \\
累积漂移 & VO 误差逐时刻传递 & \cref{ex:intro-drift} \\
尺度不确定性 & 单目差一个未知尺度因子 & \cref{thm:intro-scale} \\
单调收敛 & 路标相关性随观测单调增 & \cref{thm:intro-convergence} \\
\bottomrule
\end{tabular}
\end{table}

\begin{table}[H]\centering\small
\caption{知识点总表}
\begin{tabular}{clll}
\toprule
\# & 知识点 & 核心要点 & 难度 \\
\midrule
1 & SLAM 是什么/为什么 & 未知环境 + 携带式传感器 $\Rightarrow$ 联合估计 & $\bigstar$ \\
2 & 鸡生蛋与可解性 & 误差纠缠但相关性单调收敛 & $\bigstar\bigstar\bigstar$ \\
3 & 五模块框架 & 传感器/VO/后端/回环/建图及数据流 & $\bigstar\bigstar$ \\
4 & 累积漂移 & 纯里程计无法自纠错 & $\bigstar\bigstar$ \\
5 & 运动/观测方程 & 离散时间两条母方程 + 参数化 & $\bigstar\bigstar$ \\
6 & 状态估计两视角 & 递归贝叶斯 vs MAP/因子图 & $\bigstar\bigstar\bigstar\bigstar$ \\
7 & 滤波 vs 优化 & 线性高斯等价；优化为今主流 & $\bigstar\bigstar\bigstar$ \\
8 & 分类与代表系统 & 传感器/方法两轴 + ORB/VINS/LIO & $\bigstar\bigstar$ \\
9 & 工程基础 & 源/库/头文件/链接/可执行 & $\bigstar$ \\
\bottomrule
\end{tabular}
\end{table}

% ════════════════════════════════════════════════════════════════
\section*{延伸阅读}
\begin{itemize}
  \item \textbf{直觉与代码主骨架}：视觉 SLAM 十四讲\cite{gaoxiang2019slam14} 第 2 章——本章框架、运动/观测方程与 Hello SLAM 工程的来源。
  \item \textbf{古典时代权威导论}：Durrant-Whyte \& Bailey\cite{durrantwhyte2006slam}《SLAM: Part I》——“鸡生蛋”、单调收敛、递归贝叶斯（预测-校正）、EKF-SLAM/FastSLAM 的奠基性表述。
  \item \textbf{现代综述与未来}：Cadena 等\cite{cadena2016past}《Past, Present and Future of SLAM》——三个时代、前端/后端二分、MAP/因子图、是否需要 SLAM、是否已解决。
  \item \textbf{广度与前沿}：SLAM Handbook\cite{carlone2026handbook}——视觉 SLAM 流水线（特征法/直接法）、因子图 SLAM 的现代统一语言。
  \item \textbf{严谨主线}：Barfoot《State Estimation for Robotics》\cite{barfoot2024state}——把本章的状态估计提法严格化（概率、线性高斯、MAP、矩阵李群），是后续后端/滤波章的理论底座。
\end{itemize}

\section*{本章与后续章节的关系}
本章是全书提要，关系已在 \cref{tab:intro-roadmap} 给全；核心一句话：\textbf{框架五模块（\cref{fig:intro-pipeline}）= 后续 VO/后端/回环/建图各章的目录；数学两方程（\cref{eq:intro-slam-general}）= 数学预备（\cref{ch:rigid_body,ch:lie,ch:camera,ch:nlopt}）与状态估计（\cref{ch:slam_est,ch:eskf}）的出发点；Hello SLAM 工程 = 后续所有章的上机基础。}

\section*{故障排查手册}
\begin{enumerate}
  \item \textbf{症状}：g++ 提示 \texttt{command not found}。\textbf{原因}：未安装编译器。\textbf{对策}：\texttt{sudo apt-get install g++}（\cref{sec:intro-hello}）。
  \item \textbf{症状}：链接时报 \texttt{undefined reference to printHello()}。\textbf{原因}：忘了在 \texttt{CMakeLists.txt} 里 \texttt{target\_link\_libraries} 链接库；头文件只给声明、实现在库里。\textbf{对策}：补上链接行（\cref{sec:intro-lib} 陷阱）。
  \item \textbf{症状}：CMake 中间文件污染源码目录、难以发布。\textbf{原因}：用了 in-source 的 \texttt{cmake .}。\textbf{对策}：改用 \texttt{mkdir build; cd build; cmake ..}，发布时删 \texttt{build}（\cref{sec:intro-cmake}）。
  \item \textbf{症状}：估计轨迹整体“缩放”了、与真值差一个比例。\textbf{原因}：用单目相机，落入尺度不确定性（\cref{thm:intro-scale}）。\textbf{对策}：引入双目/RGB-D，或加 IMU 提供尺度（\cref{ch:vio}）。
  \item \textbf{症状}：长轨迹整体扭曲、直角变形、回到起点却对不上。\textbf{原因}：只有 VO、累积漂移未被校正（\cref{ex:intro-drift}）。\textbf{对策}：启用回环检测 + 后端优化（\cref{ch:loop,ch:nlopt}）。
  \item \textbf{症状}：估计结果与实际“相差十万八千里”，且越走越离谱。\textbf{原因}：把运动当成确定模型、忽略了噪声。\textbf{对策}：按随机模型携带协方差，用滤波或优化做状态估计（\cref{def:intro-state-est} 与 \cref{ch:slam_est}）。
\end{enumerate}
